\documentclass[../main.tex]{subfiles}
\begin{document}

\chapter{仮想化}
\lessoninfo{
  video={P3L6},
  textbook={OSTEP \cite{ostep} app.~B；Silberschatz ら \cite{silberschatz2018} ch.~18},
  keywords={仮想化，仮想マシン，ハイパーバイザ，ベアメタル型ハイパーバイザ，ホスト型ハイパーバイザ，保護リング，トラップアンドエミュレート，バイナリ変換，準仮想化，ハイパーコール，シャドウページテーブル，分割デバイスドライバモデル，ハードウェア支援仮想化，拡張ページテーブル},
}

これまでの章では，オペレーティングシステムが，その上で動作するプロセスに
代わってマシンの CPU，メモリ，デバイスを管理する様子を見てきた．仮想化は
この役割を1段階下へ移す．薄いソフトウェアの層が，複数の完全なオペレー
ティングシステムに代わってハードウェアを管理し，その各々がマシンを所有して
いるかのように振る舞う．この章では，仮想化を定義し，それが広く使われる
ようになった理由を説明し，仮想化ソフトウェアの2つの構成方式を比較する．
続いて，CPU，メモリ，デバイスがどのように仮想化されるかを，まずソフトウェア
だけの技法から，次に現在のプロセッサが提供する支援へと順に調べる．
第1章の保護境界，第8章のページテーブル，第11章のデバイスドライバが，
オペレーティングシステムの下にもう1層のソフトウェアを伴って，繰り返し
現れる．

\section{仮想化とは何か}
\label{sec:l12-what}

仮想化は古い考えである．1960 年代の IBM に始まる．当時，計算は少数の大型で
高価なメインフレーム上で行われており，そうしたマシンの所有者は，異なる
オペレーティングシステム向けに開発されたワークロードにそれを使いたいと
考えた．\source{P3L6，仮想化とは何か，仮想化の定義；Popek と Goldberg
\cite{popek1974}；OSTEP app.~B；Silberschatz ら ch.~18．}

\begin{definition}[仮想化]
  複数のオペレーティングシステムが，それぞれのアプリケーションとともに，
  同じ物理マシン上で並行に実行できるようにすることを
  \term[かそうか]{仮想化}[virtualization]と呼ぶ．
\end{definition}

これらのオペレーティングシステムは，もはやどれもハードウェアを制御して
いない．各々は自分専用の環境で動作し，ハードウェア資源---CPU，メモリ，
デバイス---を所有しているという幻想を与えられる．ただし，そこで見える資源は
物理資源の一部を代理する\emph{仮想資源}である（\cref{fig:l12-vmm}）．

\begin{definition}[仮想マシン]
  \term[かそうましん]{仮想マシン}[virtual machine]（VM）とは，オペレー
  ティングシステムと，そのアプリケーション，およびそれらが動作する仮想資源の
  総体である．
\end{definition}
\index{VM|see{仮想マシン}}

\begin{figure}[htbp]
  \centering
  % Virtualization: several virtual machines, each with its own guest OS,
% applications and virtual resources, on one virtual machine monitor.
% Usage: % Virtualization: several virtual machines, each with its own guest OS,
% applications and virtual resources, on one virtual machine monitor.
% Usage: % Virtualization: several virtual machines, each with its own guest OS,
% applications and virtual resources, on one virtual machine monitor.
% Usage: \input{figures/l12-vmm}   (width ~96 mm)
\begin{tikzpicture}[x=1mm, y=1mm, font=\footnotesize\sffamily,
  app/.style={draw, rounded corners=1pt, fill=white, minimum height=5mm,
              minimum width=11.5mm, inner sep=0.5pt, font=\scriptsize\sffamily},
  gos/.style={draw, fill=ln-accent!12, minimum height=6mm,
              minimum width=26mm, inner sep=0.5pt},
  vres/.style={draw, densely dashed, fill=white, minimum height=5mm,
               minimum width=26mm, inner sep=0.5pt, font=\scriptsize\sffamily},
  vm/.style={draw=ln-rule, rounded corners=2pt, fill=ln-box},
  layer/.style={draw, minimum height=7mm, minimum width=96mm, inner sep=1pt},
  ttl/.style={font=\footnotesize\sffamily\bfseries}]
  \node[layer, fill=white] at (48,3.5)
    {\iflangja{物理ハードウェア（CPU，メモリ，デバイス）}{physical hardware (CPUs, memory, devices)}};
  \node[layer, fill=ln-accent!22] at (48,12.5)
    {\iflangja{仮想マシンモニタ（ハイパーバイザ）}{virtual machine monitor (hypervisor)}};
  \foreach \x/\n in {16/1, 48/2, 80/3} {
    \draw[vm] (\x-15,18) rectangle (\x+15,48);
    \node[ttl] at (\x,44.5) {VM \n};
    \node[app] at (\x-7,37.5) {\iflangja{アプリ}{app}};
    \node[app] at (\x+7,37.5) {\iflangja{アプリ}{app}};
    \node[gos] at (\x,29.5) {\iflangja{ゲスト OS}{guest OS}};
    \node[vres] at (\x,22.5) {\iflangja{仮想資源}{virtual resources}};
  }
\end{tikzpicture}
   (width ~96 mm)
\begin{tikzpicture}[x=1mm, y=1mm, font=\footnotesize\sffamily,
  app/.style={draw, rounded corners=1pt, fill=white, minimum height=5mm,
              minimum width=11.5mm, inner sep=0.5pt, font=\scriptsize\sffamily},
  gos/.style={draw, fill=ln-accent!12, minimum height=6mm,
              minimum width=26mm, inner sep=0.5pt},
  vres/.style={draw, densely dashed, fill=white, minimum height=5mm,
               minimum width=26mm, inner sep=0.5pt, font=\scriptsize\sffamily},
  vm/.style={draw=ln-rule, rounded corners=2pt, fill=ln-box},
  layer/.style={draw, minimum height=7mm, minimum width=96mm, inner sep=1pt},
  ttl/.style={font=\footnotesize\sffamily\bfseries}]
  \node[layer, fill=white] at (48,3.5)
    {\iflangja{物理ハードウェア（CPU，メモリ，デバイス）}{physical hardware (CPUs, memory, devices)}};
  \node[layer, fill=ln-accent!22] at (48,12.5)
    {\iflangja{仮想マシンモニタ（ハイパーバイザ）}{virtual machine monitor (hypervisor)}};
  \foreach \x/\n in {16/1, 48/2, 80/3} {
    \draw[vm] (\x-15,18) rectangle (\x+15,48);
    \node[ttl] at (\x,44.5) {VM \n};
    \node[app] at (\x-7,37.5) {\iflangja{アプリ}{app}};
    \node[app] at (\x+7,37.5) {\iflangja{アプリ}{app}};
    \node[gos] at (\x,29.5) {\iflangja{ゲスト OS}{guest OS}};
    \node[vres] at (\x,22.5) {\iflangja{仮想資源}{virtual resources}};
  }
\end{tikzpicture}
   (width ~96 mm)
\begin{tikzpicture}[x=1mm, y=1mm, font=\footnotesize\sffamily,
  app/.style={draw, rounded corners=1pt, fill=white, minimum height=5mm,
              minimum width=11.5mm, inner sep=0.5pt, font=\scriptsize\sffamily},
  gos/.style={draw, fill=ln-accent!12, minimum height=6mm,
              minimum width=26mm, inner sep=0.5pt},
  vres/.style={draw, densely dashed, fill=white, minimum height=5mm,
               minimum width=26mm, inner sep=0.5pt, font=\scriptsize\sffamily},
  vm/.style={draw=ln-rule, rounded corners=2pt, fill=ln-box},
  layer/.style={draw, minimum height=7mm, minimum width=96mm, inner sep=1pt},
  ttl/.style={font=\footnotesize\sffamily\bfseries}]
  \node[layer, fill=white] at (48,3.5)
    {\iflangja{物理ハードウェア（CPU，メモリ，デバイス）}{physical hardware (CPUs, memory, devices)}};
  \node[layer, fill=ln-accent!22] at (48,12.5)
    {\iflangja{仮想マシンモニタ（ハイパーバイザ）}{virtual machine monitor (hypervisor)}};
  \foreach \x/\n in {16/1, 48/2, 80/3} {
    \draw[vm] (\x-15,18) rectangle (\x+15,48);
    \node[ttl] at (\x,44.5) {VM \n};
    \node[app] at (\x-7,37.5) {\iflangja{アプリ}{app}};
    \node[app] at (\x+7,37.5) {\iflangja{アプリ}{app}};
    \node[gos] at (\x,29.5) {\iflangja{ゲスト OS}{guest OS}};
    \node[vres] at (\x,22.5) {\iflangja{仮想資源}{virtual resources}};
  }
\end{tikzpicture}

  \caption{仮想化．各 VM は，ゲスト OS，そのアプリケーション，およびそれに
    提示される仮想資源からなる．仮想マシンモニタは，すべての VM の下で物理
    ハードウェアを管理する．}
  \label{fig:l12-vmm}
\end{figure}

\begin{definition}[仮想マシンモニタ]
  \term[かそうましんもにた]{仮想マシンモニタ}[virtual machine monitor]
  （VMM）は，\emph{ハイパーバイザ}とも呼ばれ，VM の下にある仮想化の層で
  ある．物理ハードウェアを管理し，VM に仮想資源を提供し，物理資源への
  アクセスを調停し，VM どうしを隔離する．
\end{definition}
\index{VMM|see{仮想マシンモニタ}}%
\index{はいぱーばいざ@ハイパーバイザ|see{仮想マシンモニタ}}

VM の中で動作するオペレーティングシステムを
\term[げすとOS]{ゲストOS}[guest operating system]と呼び，VM 全体は
\emph{ゲストドメイン}とも呼ばれる．

Popek と Goldberg \cite{popek1974} は古典的な定義を与えた．仮想マシンとは
「実マシンの効率的で隔離された複製」である．それを支える VMM は3つの性質を
備えなければならない．%
\margin{Popek と Goldberg はこの3つの性質を等価性，効率性，資源制御と
呼ぶ．}
\begin{description}
  \item[忠実性] プログラムは，オペレーティングシステムも含めて，元のマシンと
    本質的に同一の環境で動作し，タイミングと利用可能な資源量を除けば，元の
    マシン上でと同じように振る舞う．この\term[ちゅうじつせい]{忠実性}[fidelity]%
    があるからこそ，改変されていないオペレーティングシステムが VM の中で
    動作できる．準仮想化（\cref{sec:l12-pv}）はこれを意図的に緩める．
  \item[性能] VM 内のプログラムの速度低下は，最悪でもわずかである．したがって
    ほとんどの命令は，VMM が解釈するのではなく，ハードウェア上で直接実行され
    なければならない．
  \item[安全性と隔離] VMM がシステム資源を完全に制御する．VM は，VMM が割り
    当てていない資源にアクセスできず，他の VM や VMM 自体に干渉することも
    できない．
\end{description}

\begin{note}
  この定義は，同じく仮想マシンと呼ばれる技術を除外する．高水準言語仮想マシン
  は，実マシン向けではなく抽象的な命令セット向けにコンパイルされたプログラム
  を実行するので，その下のハードウェアの複製ではない．異なるハードウェア
  プラットフォームをソフトウェアで実装するエミュレータは，実マシンではなく
  そのプラットフォームを複製しており，しかもすべての命令を解釈または変換に
  よってソフトウェアで実行し，どれ1つとしてプロセッサ上で直接実行しないので，
  性能の要件を満たさない．
\end{note}

\section{仮想化の利点}
\label{sec:l12-benefits}

1台の物理マシン上で複数のオペレーティングシステムを動かすことには，いくつもの
利点がある．\source{P3L6，仮想化の利点；Rosenblum と Garfinkel
\cite{rosenblum2005}；Silberschatz ら ch.~18．}
\begin{description}
  \item[サーバ統合] それぞれ独自のオペレーティングシステムとアプリケーションを
    持つ複数の物理マシンのワークロードを，より少ない台数のマシン上で VM として
    動作させられる．この\term[さーばとうごう]{サーバ統合}[consolidation]は
    コストを下げ---購入，設置，給電，冷却すべきマシンが減る---，管理すべき
    マシンが減るので管理性も高まる．
  \item[マイグレーション] VM は物理マシンから切り離されているので，オペレー
    ティングシステムとアプリケーションを含むその状態全体を，例えば保守のため
    にマシンを停止する前に，あるマシンから別のマシンへ移せる．マイグレーション
    は可用性と信頼性を高める．
  \item[セキュリティ] バグや悪意ある振る舞いは，それが起こった VM の中に
    封じ込められ，他の VM に影響しない．
  \item[デバッグ] 新しいオペレーティングシステムやその新しい機能を，専用の
    マシンなしに VM 上で素早く立ち上げて試験できる．クラッシュしても落ちるのは
    その VM だけである．
  \item[レガシー OS のサポート] 現在のハードウェア上ではもはや動作しない
    オペレーティングシステムを，それに依存するアプリケーションとともに，現行の
    システムの隣で VM の中で動かし続けられる．
\end{description}

IBM のメインフレームは 1960 年代から仮想化をサポートしていたが，1990 年代
後半までそれはメインフレームに限られていた
\cite{rosenblum2005}．\sidenote{最初のそのようなシステムは CP-40 と CMS で
あり，IBM ケンブリッジ科学センターが 1967 年 1 月から運用した．続く CP-67
は，ハードウェアがアドレスを変換する System/360 モデル 67 の上で動いた．この
系譜は 1972 年に発表された VM/370 へ続き，現在の z/VM に至る．}メイン
フレームは高価で数も少なかった．コモディティハードウェアが安価になり，その
オペレーティングシステムがマルチタスクをサポートするようになると，1台の
マシンを複数のオペレーティングシステムで分け合うよりも，別のアプリケーション
のために別のマシンを買う方が簡単であり，VMM への関心は薄れた．この慣行が，
やがて VMM を復活させる問題を生んだ．データセンタは，それぞれ1つの
アプリケーションだけを動かす低稼働率のマシンで埋まり，それらが消費する場所，
電力，管理の手間が，果たす仕事に見合わないほど大きくなった．アプリケーション
をより少ないマシンに統合すればこれらのコストは一挙に減るが，そのためには
コモディティな x86 ハードウェアを仮想化しなければならず，\cref{sec:l12-x86}で
見るように，x86 はそのために設計されてはいなかった．

\begin{example}
  ある企業が，それぞれ1つのアプリケーション専用の 24 台のサーバを運用して
  おり，CPU 使用率は平均 \SI{8}{\percent} である．\margin{低い稼働率は例外では
  ない．Barroso と H\"olzle は，データセンタのサーバが時間の大半を CPU
  使用率 \SI{10}{\percent} から \SI{50}{\percent} の間で過ごすと報告して
  いる（\emph{The Datacenter as a Computer}）．}合計の負荷は
  $24 \times 0.08 = 1.92$ 台分の全力稼働に相当する．負荷のピークと仮想化の
  オーバーヘッドに余裕を残して各物理マシンの負荷を高々 \SI{70}{\percent} に
  抑え，メモリと I/O の容量が足りるなら，24 のワークロードは
  $\lceil 1.92/0.7 \rceil = 3$ 台のマシンに VM として収まり，各マシンの
  使用率は約 \SI{64}{\percent} となる．給電，冷却，管理すべきマシンの台数は
  8 分の 1 に減る．
\end{example}

\section{仮想化のモデル}
\label{sec:l12-models}

仮想化ソフトウェアの構成には2つの方式があり，何がハードウェア上で直接動作
するか，とりわけ誰がデバイスドライバを所有するかが異なる．\source{P3L6，
仮想化のモデル: ベアメタル型，ホスト型；Barham ら \cite{barham2003}；
Kivity ら \cite{kivity2007}；Silberschatz ら ch.~18．}

\subsection{ベアメタル型ハイパーバイザ}

\begin{definition}[ベアメタル型ハイパーバイザ]
  \term[べあめたるがたはいぱーばいざ]{ベアメタル型ハイパーバイザ}%
  [bare-metal hypervisor]は，ハイパーバイザ型 VMM あるいは\emph{タイプ 1}
  VMM とも呼ばれ，ハードウェア上で直接動作する．すべてのハードウェア資源を
  管理し，VM の実行を支える．
\end{definition}

このようなハイパーバイザはデバイスも制御しなければならない．デバイスは多種
多様であり，そのすべてをサポートすれば大量のデバイスドライバがハイパーバイザ
に入り込む．そこでベアメタル型の設計では，ゲスト VM の隣に特権的な
\emph{サービス VM} を動かすことが多い（\cref{fig:l12-models}a）．サービス VM
は，デバイスドライバを備えた標準的なオペレーティングシステムを動かし，
デバイスに直接アクセスでき，構成と管理の作業を担う．こうしてハイパーバイザ
自体は小さく保たれる．\caution{タイプ 1 とタイプ 2 という呼び分けは，きれいに
割り切れるものではない．KVM は Linux のモジュールなので，そのハイパーバイザは
ハードウェア上で動作している．それでもタイプ 2 に数えられるのは，その周りに
完全なホスト OS が残っているからである．}

\begin{figure}[htbp]
  \centering
  % Virtualization models: (a) bare-metal hypervisor with a privileged
% service VM that owns the device drivers (Xen: dom0 and domU guests);
% (b) hosted hypervisor: a module in the host OS (KVM in Linux), each VM a
% host process with QEMU emulating its devices.
% Usage: % Virtualization models: (a) bare-metal hypervisor with a privileged
% service VM that owns the device drivers (Xen: dom0 and domU guests);
% (b) hosted hypervisor: a module in the host OS (KVM in Linux), each VM a
% host process with QEMU emulating its devices.
% Usage: % Virtualization models: (a) bare-metal hypervisor with a privileged
% service VM that owns the device drivers (Xen: dom0 and domU guests);
% (b) hosted hypervisor: a module in the host OS (KVM in Linux), each VM a
% host process with QEMU emulating its devices.
% Usage: \input{figures/l12-models}   (width ~116 mm)
\begin{tikzpicture}[x=1mm, y=1mm, font=\footnotesize\sffamily,
  >={Stealth[length=1.6mm]},
  box/.style={draw, fill=white, inner sep=0.5pt, align=center,
              font=\scriptsize\sffamily},
  gos/.style={box, fill=ln-accent!12},
  vm/.style={draw=ln-rule, rounded corners=2pt, fill=ln-box},
  svm/.style={draw=ln-accent, rounded corners=2pt, fill=ln-accent!8},
  lab/.style={font=\scriptsize\sffamily, align=center},
  hdr/.style={font=\scriptsize\sffamily\bfseries, align=center},
  pcap/.style={font=\footnotesize\sffamily}]
  % ================= (a) bare-metal
  \draw[fill=white] (0,0) rectangle (55,7)
    node[midway] {\iflangja{ハードウェア}{hardware}};
  \draw[fill=ln-accent!22] (0,9) rectangle (55,16)
    node[midway] {\iflangja{ハイパーバイザ（VMM）}{hypervisor (VMM)}};
  % service VM
  \draw[svm] (0,18) rectangle (19,52);
  \node[hdr] at (9.5,49) {\iflangja{サービス VM}{service VM}};
  \node[lab] at (9.5,45.5) {(dom0)};
  \node[box, minimum width=16mm, minimum height=7mm] at (9.5,36.5)
    {\iflangja{管理ツール}{management}};
  \node[gos, minimum width=16mm, minimum height=9mm] (drv) at (9.5,25.5)
    {OS +\\\iflangja{ドライバ}{drivers}};
  \draw[<->, thick, densely dashed, ln-accent] (drv.south) -- (9.5,7);
  % guest VMs
  \foreach \x in {29, 47} {
    \draw[vm] (\x-8,18) rectangle (\x+8,52);
    \node[hdr] at (\x,49) {\iflangja{ゲスト VM}{guest VM}};
    \node[lab] at (\x,45.5) {(domU)};
    \node[box, rounded corners=1pt, minimum width=13mm, minimum height=6mm]
      at (\x,36.5) {\iflangja{アプリ}{apps}};
    \node[gos, minimum width=13mm, minimum height=7mm] at (\x,26)
      {\iflangja{ゲスト OS}{guest OS}};
  }
  \node[pcap] at (27.5,-4) {\iflangja{(a) ベアメタル型}{(a) bare-metal}};
  % ================= (b) hosted
  \begin{scope}[xshift=61mm]
  \draw[fill=white] (0,0) rectangle (55,7)
    node[midway] {\iflangja{ハードウェア}{hardware}};
  \draw[fill=ln-box] (0,9) rectangle (55,26);
  \node[hdr] at (27.5,22.5) {\iflangja{ホスト OS（Linux）}{host OS (Linux)}};
  \node[box, minimum width=24mm, minimum height=7mm] at (14,14.5)
    {\iflangja{デバイスドライバ}{device drivers}};
  \node[box, fill=ln-accent!22, minimum width=24mm, minimum height=7mm]
    at (41,14.5) {\iflangja{KVM モジュール}{KVM module}};
  % native process
  \draw[fill=white, rounded corners=2pt] (0,28) rectangle (13,52)
    node[midway, lab] {\iflangja{通常の\\プロセス}{native\\process}};
  % VMs as host processes
  \foreach \x in {24.5, 45.5} {
    \draw[vm] (\x-9.5,28) rectangle (\x+9.5,52);
    \node[hdr] at (\x,49.3) {VM};
    \node[box, rounded corners=1pt, minimum width=15mm, minimum height=5mm]
      at (\x,43.5) {\iflangja{アプリ}{apps}};
    \node[gos, minimum width=16mm, minimum height=5.5mm] at (\x,37)
      {\iflangja{ゲスト OS}{guest OS}};
    \node[box, densely dashed, minimum width=16mm, minimum height=5mm]
      at (\x,31) {QEMU};
  }
  \node[pcap] at (27.5,-4) {\iflangja{(b) ホスト型}{(b) hosted}};
  \end{scope}
\end{tikzpicture}
   (width ~116 mm)
\begin{tikzpicture}[x=1mm, y=1mm, font=\footnotesize\sffamily,
  >={Stealth[length=1.6mm]},
  box/.style={draw, fill=white, inner sep=0.5pt, align=center,
              font=\scriptsize\sffamily},
  gos/.style={box, fill=ln-accent!12},
  vm/.style={draw=ln-rule, rounded corners=2pt, fill=ln-box},
  svm/.style={draw=ln-accent, rounded corners=2pt, fill=ln-accent!8},
  lab/.style={font=\scriptsize\sffamily, align=center},
  hdr/.style={font=\scriptsize\sffamily\bfseries, align=center},
  pcap/.style={font=\footnotesize\sffamily}]
  % ================= (a) bare-metal
  \draw[fill=white] (0,0) rectangle (55,7)
    node[midway] {\iflangja{ハードウェア}{hardware}};
  \draw[fill=ln-accent!22] (0,9) rectangle (55,16)
    node[midway] {\iflangja{ハイパーバイザ（VMM）}{hypervisor (VMM)}};
  % service VM
  \draw[svm] (0,18) rectangle (19,52);
  \node[hdr] at (9.5,49) {\iflangja{サービス VM}{service VM}};
  \node[lab] at (9.5,45.5) {(dom0)};
  \node[box, minimum width=16mm, minimum height=7mm] at (9.5,36.5)
    {\iflangja{管理ツール}{management}};
  \node[gos, minimum width=16mm, minimum height=9mm] (drv) at (9.5,25.5)
    {OS +\\\iflangja{ドライバ}{drivers}};
  \draw[<->, thick, densely dashed, ln-accent] (drv.south) -- (9.5,7);
  % guest VMs
  \foreach \x in {29, 47} {
    \draw[vm] (\x-8,18) rectangle (\x+8,52);
    \node[hdr] at (\x,49) {\iflangja{ゲスト VM}{guest VM}};
    \node[lab] at (\x,45.5) {(domU)};
    \node[box, rounded corners=1pt, minimum width=13mm, minimum height=6mm]
      at (\x,36.5) {\iflangja{アプリ}{apps}};
    \node[gos, minimum width=13mm, minimum height=7mm] at (\x,26)
      {\iflangja{ゲスト OS}{guest OS}};
  }
  \node[pcap] at (27.5,-4) {\iflangja{(a) ベアメタル型}{(a) bare-metal}};
  % ================= (b) hosted
  \begin{scope}[xshift=61mm]
  \draw[fill=white] (0,0) rectangle (55,7)
    node[midway] {\iflangja{ハードウェア}{hardware}};
  \draw[fill=ln-box] (0,9) rectangle (55,26);
  \node[hdr] at (27.5,22.5) {\iflangja{ホスト OS（Linux）}{host OS (Linux)}};
  \node[box, minimum width=24mm, minimum height=7mm] at (14,14.5)
    {\iflangja{デバイスドライバ}{device drivers}};
  \node[box, fill=ln-accent!22, minimum width=24mm, minimum height=7mm]
    at (41,14.5) {\iflangja{KVM モジュール}{KVM module}};
  % native process
  \draw[fill=white, rounded corners=2pt] (0,28) rectangle (13,52)
    node[midway, lab] {\iflangja{通常の\\プロセス}{native\\process}};
  % VMs as host processes
  \foreach \x in {24.5, 45.5} {
    \draw[vm] (\x-9.5,28) rectangle (\x+9.5,52);
    \node[hdr] at (\x,49.3) {VM};
    \node[box, rounded corners=1pt, minimum width=15mm, minimum height=5mm]
      at (\x,43.5) {\iflangja{アプリ}{apps}};
    \node[gos, minimum width=16mm, minimum height=5.5mm] at (\x,37)
      {\iflangja{ゲスト OS}{guest OS}};
    \node[box, densely dashed, minimum width=16mm, minimum height=5mm]
      at (\x,31) {QEMU};
  }
  \node[pcap] at (27.5,-4) {\iflangja{(b) ホスト型}{(b) hosted}};
  \end{scope}
\end{tikzpicture}
   (width ~116 mm)
\begin{tikzpicture}[x=1mm, y=1mm, font=\footnotesize\sffamily,
  >={Stealth[length=1.6mm]},
  box/.style={draw, fill=white, inner sep=0.5pt, align=center,
              font=\scriptsize\sffamily},
  gos/.style={box, fill=ln-accent!12},
  vm/.style={draw=ln-rule, rounded corners=2pt, fill=ln-box},
  svm/.style={draw=ln-accent, rounded corners=2pt, fill=ln-accent!8},
  lab/.style={font=\scriptsize\sffamily, align=center},
  hdr/.style={font=\scriptsize\sffamily\bfseries, align=center},
  pcap/.style={font=\footnotesize\sffamily}]
  % ================= (a) bare-metal
  \draw[fill=white] (0,0) rectangle (55,7)
    node[midway] {\iflangja{ハードウェア}{hardware}};
  \draw[fill=ln-accent!22] (0,9) rectangle (55,16)
    node[midway] {\iflangja{ハイパーバイザ（VMM）}{hypervisor (VMM)}};
  % service VM
  \draw[svm] (0,18) rectangle (19,52);
  \node[hdr] at (9.5,49) {\iflangja{サービス VM}{service VM}};
  \node[lab] at (9.5,45.5) {(dom0)};
  \node[box, minimum width=16mm, minimum height=7mm] at (9.5,36.5)
    {\iflangja{管理ツール}{management}};
  \node[gos, minimum width=16mm, minimum height=9mm] (drv) at (9.5,25.5)
    {OS +\\\iflangja{ドライバ}{drivers}};
  \draw[<->, thick, densely dashed, ln-accent] (drv.south) -- (9.5,7);
  % guest VMs
  \foreach \x in {29, 47} {
    \draw[vm] (\x-8,18) rectangle (\x+8,52);
    \node[hdr] at (\x,49) {\iflangja{ゲスト VM}{guest VM}};
    \node[lab] at (\x,45.5) {(domU)};
    \node[box, rounded corners=1pt, minimum width=13mm, minimum height=6mm]
      at (\x,36.5) {\iflangja{アプリ}{apps}};
    \node[gos, minimum width=13mm, minimum height=7mm] at (\x,26)
      {\iflangja{ゲスト OS}{guest OS}};
  }
  \node[pcap] at (27.5,-4) {\iflangja{(a) ベアメタル型}{(a) bare-metal}};
  % ================= (b) hosted
  \begin{scope}[xshift=61mm]
  \draw[fill=white] (0,0) rectangle (55,7)
    node[midway] {\iflangja{ハードウェア}{hardware}};
  \draw[fill=ln-box] (0,9) rectangle (55,26);
  \node[hdr] at (27.5,22.5) {\iflangja{ホスト OS（Linux）}{host OS (Linux)}};
  \node[box, minimum width=24mm, minimum height=7mm] at (14,14.5)
    {\iflangja{デバイスドライバ}{device drivers}};
  \node[box, fill=ln-accent!22, minimum width=24mm, minimum height=7mm]
    at (41,14.5) {\iflangja{KVM モジュール}{KVM module}};
  % native process
  \draw[fill=white, rounded corners=2pt] (0,28) rectangle (13,52)
    node[midway, lab] {\iflangja{通常の\\プロセス}{native\\process}};
  % VMs as host processes
  \foreach \x in {24.5, 45.5} {
    \draw[vm] (\x-9.5,28) rectangle (\x+9.5,52);
    \node[hdr] at (\x,49.3) {VM};
    \node[box, rounded corners=1pt, minimum width=15mm, minimum height=5mm]
      at (\x,43.5) {\iflangja{アプリ}{apps}};
    \node[gos, minimum width=16mm, minimum height=5.5mm] at (\x,37)
      {\iflangja{ゲスト OS}{guest OS}};
    \node[box, densely dashed, minimum width=16mm, minimum height=5mm]
      at (\x,31) {QEMU};
  }
  \node[pcap] at (27.5,-4) {\iflangja{(b) ホスト型}{(b) hosted}};
  \end{scope}
\end{tikzpicture}

  \caption{仮想化のモデル．(a) デバイスドライバを所有する特権的なサービス VM
    を伴うベアメタル型ハイパーバイザ．Xen がこの形である．(b) ホスト型
    ハイパーバイザ: Linux における KVM のようにホスト OS のモジュールであり，
    各 VM はホストのプロセスで，その中で QEMU がデバイスをエミュレートする．}
  \label{fig:l12-models}
\end{figure}

ハイパーバイザ \term{Xen}[Xen] \cite{barham2003} は，オープンソースとしても
商用の Citrix XenServer としても入手でき，この設計に従う．その VM は
\emph{ドメイン}と呼ばれる．特権的なサービス VM は \term{dom0}[dom0] であり，
デバイスドライバと管理ツールを備えた Linux を動かす．ゲスト VM は非特権の
ドメイン \emph{domU} であり，Xen ハイパーバイザ自体は，ディスクやネットワーク
カードといった I/O デバイスのドライバを一切含まない．\margin{Barham らが
記述した最初の Xen \cite{barham2003} では，これらのドライバはまだ
ハイパーバイザの中にあった．}VMware ESX はもう一方の道を採り，ドライバを
ハイパーバイザに置く．ESX が対象とするサーバのベンダは限られており，各ベンダ
が ESX 向けのドライバを供給するので，これが実用的である．ESX は遠隔管理 API
を通じて管理され，この API が以前の Linux ベースの制御用中核を置き換えた．

\subsection{ホスト型ハイパーバイザ}

\begin{definition}[ホスト型ハイパーバイザ]
  \term[ほすとがたはいぱーばいざ]{ホスト型ハイパーバイザ}[hosted hypervisor]，
  すなわち\emph{タイプ 2} VMM は，従来型のオペレーティングシステムの上で動作
  する．このオペレーティングシステムを
  \term[ほすとOS]{ホストOS}[host operating system]と呼び，ホスト OS が
  ハードウェアのすべてを所有する．ホスト OS 内の特別な VMM モジュールが，VM に
  ハードウェアインタフェースを提供し，VM 間の切り替えを扱う．
\end{definition}

ホスト型ハイパーバイザは，ホスト OS がすでに提供しているものを再利用する．
デバイスドライバ，スケジューリング，メモリ管理などのサービスを実装し直す必要
はなく，通常のアプリケーションも VM の隣でホスト上でそのまま動き続ける
（\cref{fig:l12-models}b）．

\needspace{6\baselineskip}
\term{KVM}[Kernel-based Virtual Machine]%
\index{Kernel-based Virtual Machine|see{KVM}} \cite{kivity2007} は，Linux の
ホスト型ハイパーバイザである．ロード可能なカーネルモジュールが Linux を
ハイパーバイザに変える．各 VM は通常の Linux プロセスであり，その仮想 CPU は
そのプロセスのスレッドである．したがって Linux のスケジューラとメモリ
マネージャは，他のプロセスと同じように VM を扱う．VM のデバイスは，ユーザ
レベルで動作するハードウェアエミュレータ \term{QEMU}[QEMU] によってその
プロセスの中でエミュレートされ，一方ゲストの命令は\cref{sec:l12-hw}で述べる
ハードウェア支援によって CPU 上で直接実行される．KVM は Linux の一部である
ため，Linux のオープンソースコミュニティの恩恵を受ける．新しい機能，バグ修正，
新しいデバイスのドライバが，別途の労力なしに利用できるようになる．

\begin{keypoint}
  ベアメタル型ハイパーバイザはハードウェア上で直接動作し，デバイスドライバは
  Xen の dom0 のような特権的なサービス VM に委ねるのが普通である．ホスト型
  ハイパーバイザは Linux における KVM のようにホスト OS のモジュールであり，
  その OS のドライバとサービスを再利用する．
\end{keypoint}

\section{ハードウェアの保護レベル}
\label{sec:l12-rings}

ゲスト OS は，カーネルモードで動作し特権命令を実行するように書かれているが，
仮想化の下ではハイパーバイザだけがハードウェアを制御してよい．\source{P3L6，
ハードウェアの保護レベル；Silberschatz ら ch.~18．}第1章の2つのモードでは，
アプリケーション，ゲスト OS，ハイパーバイザという3種類のソフトウェアを分離
できないが，コモディティプロセッサは2つより多くの特権レベルを備えている．

\begin{definition}[保護リング]
  \term[ほごりんぐ]{保護リング}[protection ring]とは，プロセッサが強制する
  特権レベルの階層の1つであり，あるリングのソフトウェアは，そのレベルで許され
  た命令だけを実行し，許された資源だけにアクセスできる．x86 プロセッサは4つの
  リングを備え，最も特権的なリング 0 から最も特権の低いリング 3 まである．
\end{definition}

仮想化を用いない場合，OS カーネルはリング 0，アプリケーションはリング 3 で
動作し，リング 1 と 2 は使われない（\cref{fig:l12-rings}a）．仮想化を用いる
場合，ハイパーバイザがリング 0 を占め，ゲスト OS はリング 1 へ移り，
アプリケーションはリング 3 のままである（\cref{fig:l12-rings}b）．ゲスト OS が
リング 1 から特権的な操作を試みると，プロセッサはリング 0 のハイパーバイザへ
トラップする．

\begin{figure}[htbp]
  \centering
  % x86 protection rings: (a) native OS; (b) virtualization without hardware
% support (hypervisor in ring 0, guest OS deprivileged to ring 1);
% (c) root and non-root modes, each with four rings, VM exit / VM entry.
% Usage: % x86 protection rings: (a) native OS; (b) virtualization without hardware
% support (hypervisor in ring 0, guest OS deprivileged to ring 1);
% (c) root and non-root modes, each with four rings, VM exit / VM entry.
% Usage: % x86 protection rings: (a) native OS; (b) virtualization without hardware
% support (hypervisor in ring 0, guest OS deprivileged to ring 1);
% (c) root and non-root modes, each with four rings, VM exit / VM entry.
% Usage: \input{figures/l12-rings}   (width ~118 mm)
\begin{tikzpicture}[x=1mm, y=1mm, font=\scriptsize\sffamily,
  >={Stealth[length=1.5mm]},
  used/.style={draw, fill=ln-accent!12},
  hv/.style={draw, fill=ln-accent!30},
  free/.style={draw=ln-rule, fill=white},
  rl/.style={font=\scriptsize\sffamily, anchor=east, text=ln-muted},
  hdr/.style={font=\scriptsize\sffamily\bfseries},
  pcap/.style={font=\footnotesize\sffamily, align=center, anchor=north}]
  % row labels
  \foreach \r/\y in {3/27, 2/21, 1/15, 0/9}
    \node[rl] at (9.5,\y) {\iflangja{リング \r}{ring \r}};
  % (a) native
  \draw[used] (11,24) rectangle (31,30) node[midway] {\iflangja{アプリケーション}{applications}};
  \draw[free] (11,18) rectangle (31,24);
  \draw[free] (11,12) rectangle (31,18);
  \draw[used] (11,6)  rectangle (31,12) node[midway] {\iflangja{OS カーネル}{OS kernel}};
  \node[pcap] at (21,-1) {\iflangja{(a) 仮想化なし}{(a) native}};
  % (b) hypervisor in ring 0
  \draw[used] (34,24) rectangle (54,30) node[midway] {\iflangja{アプリケーション}{applications}};
  \draw[free] (34,18) rectangle (54,24);
  \draw[used] (34,12) rectangle (54,18) node[midway] {\iflangja{ゲスト OS}{guest OS}};
  \draw[hv]   (34,6)  rectangle (54,12) node[midway] {\iflangja{ハイパーバイザ}{hypervisor}};
  \draw[->, thick, ln-caution] (54.3,15) to[out=0, in=0, looseness=2.2]
    node[right=0.8mm, text=ln-caution] {\iflangja{}{trap}} (54.3,9);
  % The Japanese label is wider than the gap between panels (b) and (c),
  % so it is set vertically beside the arc instead of next to it.
  \iflangja{\node[text=ln-caution, rotate=90] at (60.3,12) {トラップ};}{}
  \node[pcap] at (44,-1) {\iflangja{(b) リング 0 の\\ハイパーバイザ}{(b) hypervisor\\in ring 0}};
  % (c) root / non-root
  \node[hdr] at (74,33) {\iflangja{非ルートモード}{non-root mode}};
  \node[hdr] at (108,33) {\iflangja{ルートモード}{root mode}};
  \draw[used] (64,24) rectangle (84,30) node[midway] {\iflangja{アプリケーション}{applications}};
  \draw[free] (64,18) rectangle (84,24);
  \draw[free] (64,12) rectangle (84,18);
  \draw[used] (64,6)  rectangle (84,12) node[midway] {\iflangja{ゲスト OS}{guest OS}};
  \draw[free] (98,24) rectangle (118,30);
  \draw[free] (98,18) rectangle (118,24);
  \draw[free] (98,12) rectangle (118,18);
  \draw[hv]   (98,6)  rectangle (118,12) node[midway] {\iflangja{ハイパーバイザ}{hypervisor}};
  \draw[->, thick, ln-caution] (84.3,10.2) -- (97.7,10.2);
  \draw[->, thick, ln-tip] (97.7,7.8) -- (84.3,7.8);
  \node[text=ln-caution, anchor=south] at (91,12.2) {VM exit};
  \node[text=ln-tip, anchor=north] at (91,5.8) {VM entry};
  \node[pcap] at (91,-1) {\iflangja{(c) ルートモードと非ルートモード}{(c) root and non-root modes}};
\end{tikzpicture}
   (width ~118 mm)
\begin{tikzpicture}[x=1mm, y=1mm, font=\scriptsize\sffamily,
  >={Stealth[length=1.5mm]},
  used/.style={draw, fill=ln-accent!12},
  hv/.style={draw, fill=ln-accent!30},
  free/.style={draw=ln-rule, fill=white},
  rl/.style={font=\scriptsize\sffamily, anchor=east, text=ln-muted},
  hdr/.style={font=\scriptsize\sffamily\bfseries},
  pcap/.style={font=\footnotesize\sffamily, align=center, anchor=north}]
  % row labels
  \foreach \r/\y in {3/27, 2/21, 1/15, 0/9}
    \node[rl] at (9.5,\y) {\iflangja{リング \r}{ring \r}};
  % (a) native
  \draw[used] (11,24) rectangle (31,30) node[midway] {\iflangja{アプリケーション}{applications}};
  \draw[free] (11,18) rectangle (31,24);
  \draw[free] (11,12) rectangle (31,18);
  \draw[used] (11,6)  rectangle (31,12) node[midway] {\iflangja{OS カーネル}{OS kernel}};
  \node[pcap] at (21,-1) {\iflangja{(a) 仮想化なし}{(a) native}};
  % (b) hypervisor in ring 0
  \draw[used] (34,24) rectangle (54,30) node[midway] {\iflangja{アプリケーション}{applications}};
  \draw[free] (34,18) rectangle (54,24);
  \draw[used] (34,12) rectangle (54,18) node[midway] {\iflangja{ゲスト OS}{guest OS}};
  \draw[hv]   (34,6)  rectangle (54,12) node[midway] {\iflangja{ハイパーバイザ}{hypervisor}};
  \draw[->, thick, ln-caution] (54.3,15) to[out=0, in=0, looseness=2.2]
    node[right=0.8mm, text=ln-caution] {\iflangja{}{trap}} (54.3,9);
  % The Japanese label is wider than the gap between panels (b) and (c),
  % so it is set vertically beside the arc instead of next to it.
  \iflangja{\node[text=ln-caution, rotate=90] at (60.3,12) {トラップ};}{}
  \node[pcap] at (44,-1) {\iflangja{(b) リング 0 の\\ハイパーバイザ}{(b) hypervisor\\in ring 0}};
  % (c) root / non-root
  \node[hdr] at (74,33) {\iflangja{非ルートモード}{non-root mode}};
  \node[hdr] at (108,33) {\iflangja{ルートモード}{root mode}};
  \draw[used] (64,24) rectangle (84,30) node[midway] {\iflangja{アプリケーション}{applications}};
  \draw[free] (64,18) rectangle (84,24);
  \draw[free] (64,12) rectangle (84,18);
  \draw[used] (64,6)  rectangle (84,12) node[midway] {\iflangja{ゲスト OS}{guest OS}};
  \draw[free] (98,24) rectangle (118,30);
  \draw[free] (98,18) rectangle (118,24);
  \draw[free] (98,12) rectangle (118,18);
  \draw[hv]   (98,6)  rectangle (118,12) node[midway] {\iflangja{ハイパーバイザ}{hypervisor}};
  \draw[->, thick, ln-caution] (84.3,10.2) -- (97.7,10.2);
  \draw[->, thick, ln-tip] (97.7,7.8) -- (84.3,7.8);
  \node[text=ln-caution, anchor=south] at (91,12.2) {VM exit};
  \node[text=ln-tip, anchor=north] at (91,5.8) {VM entry};
  \node[pcap] at (91,-1) {\iflangja{(c) ルートモードと非ルートモード}{(c) root and non-root modes}};
\end{tikzpicture}
   (width ~118 mm)
\begin{tikzpicture}[x=1mm, y=1mm, font=\scriptsize\sffamily,
  >={Stealth[length=1.5mm]},
  used/.style={draw, fill=ln-accent!12},
  hv/.style={draw, fill=ln-accent!30},
  free/.style={draw=ln-rule, fill=white},
  rl/.style={font=\scriptsize\sffamily, anchor=east, text=ln-muted},
  hdr/.style={font=\scriptsize\sffamily\bfseries},
  pcap/.style={font=\footnotesize\sffamily, align=center, anchor=north}]
  % row labels
  \foreach \r/\y in {3/27, 2/21, 1/15, 0/9}
    \node[rl] at (9.5,\y) {\iflangja{リング \r}{ring \r}};
  % (a) native
  \draw[used] (11,24) rectangle (31,30) node[midway] {\iflangja{アプリケーション}{applications}};
  \draw[free] (11,18) rectangle (31,24);
  \draw[free] (11,12) rectangle (31,18);
  \draw[used] (11,6)  rectangle (31,12) node[midway] {\iflangja{OS カーネル}{OS kernel}};
  \node[pcap] at (21,-1) {\iflangja{(a) 仮想化なし}{(a) native}};
  % (b) hypervisor in ring 0
  \draw[used] (34,24) rectangle (54,30) node[midway] {\iflangja{アプリケーション}{applications}};
  \draw[free] (34,18) rectangle (54,24);
  \draw[used] (34,12) rectangle (54,18) node[midway] {\iflangja{ゲスト OS}{guest OS}};
  \draw[hv]   (34,6)  rectangle (54,12) node[midway] {\iflangja{ハイパーバイザ}{hypervisor}};
  \draw[->, thick, ln-caution] (54.3,15) to[out=0, in=0, looseness=2.2]
    node[right=0.8mm, text=ln-caution] {\iflangja{}{trap}} (54.3,9);
  % The Japanese label is wider than the gap between panels (b) and (c),
  % so it is set vertically beside the arc instead of next to it.
  \iflangja{\node[text=ln-caution, rotate=90] at (60.3,12) {トラップ};}{}
  \node[pcap] at (44,-1) {\iflangja{(b) リング 0 の\\ハイパーバイザ}{(b) hypervisor\\in ring 0}};
  % (c) root / non-root
  \node[hdr] at (74,33) {\iflangja{非ルートモード}{non-root mode}};
  \node[hdr] at (108,33) {\iflangja{ルートモード}{root mode}};
  \draw[used] (64,24) rectangle (84,30) node[midway] {\iflangja{アプリケーション}{applications}};
  \draw[free] (64,18) rectangle (84,24);
  \draw[free] (64,12) rectangle (84,18);
  \draw[used] (64,6)  rectangle (84,12) node[midway] {\iflangja{ゲスト OS}{guest OS}};
  \draw[free] (98,24) rectangle (118,30);
  \draw[free] (98,18) rectangle (118,24);
  \draw[free] (98,12) rectangle (118,18);
  \draw[hv]   (98,6)  rectangle (118,12) node[midway] {\iflangja{ハイパーバイザ}{hypervisor}};
  \draw[->, thick, ln-caution] (84.3,10.2) -- (97.7,10.2);
  \draw[->, thick, ln-tip] (97.7,7.8) -- (84.3,7.8);
  \node[text=ln-caution, anchor=south] at (91,12.2) {VM exit};
  \node[text=ln-tip, anchor=north] at (91,5.8) {VM entry};
  \node[pcap] at (91,-1) {\iflangja{(c) ルートモードと非ルートモード}{(c) root and non-root modes}};
\end{tikzpicture}

  \caption{x86 の保護リングの使い方．(a) 仮想化なしの実行．
    (b) ハードウェア支援のない仮想化: ゲスト OS はリング 1 へ降格され，その
    特権的な操作はハイパーバイザへトラップする．(c) 仮想化支援を備えた
    プロセッサ: ゲスト OS は非ルートモードのリング 0 で動作し，VM exit が制御を
    ルートモードのハイパーバイザへ移す．}
  \label{fig:l12-rings}
\end{figure}

仮想化支援を備えたプロセッサは，現在の x86-64 プロセッサを含め，さらに進んで，
それぞれ4つのリングを持つ2つの動作モードを提供する（\cref{fig:l12-rings}c）．

\begin{definition}[ルートモードと非ルートモード]
  \term[るーともーど]{ルートモード}[root mode]では，プロセッサはハイパー
  バイザをリング 0 で，ハードウェアを完全に制御した状態で実行する．
  \term[ひるーともーど]{非ルートモード}[non-root mode]では VM を実行し，
  ゲスト OS はリング 0，アプリケーションはリング 3 で動作するが，特権的な操作は
  制限される．
\end{definition}

こうしてゲスト OS は，それが書かれたときに想定された特権レベルで動作する．
ゲスト OS が特権的な操作を行おうとすると，\term{VM exit}[VM exit] と呼ばれる
トラップが起こり，プロセッサはルートモードへ切り替わって制御がハイパーバイザ
に渡る．ハイパーバイザがその操作を処理し終えると，
\emph{VM entry}\index{VM entry} がプロセッサを非ルートモードへ戻し，VM を
再開する．

\section{プロセッサの仮想化}
\label{sec:l12-cpu}

\cref{sec:l12-what}の性能の要件を満たすため，ハイパーバイザは VM の命令を
解釈するのではなく，ハードウェア上で直接実行させる．\source{P3L6，プロセッサ
の仮想化；Popek と Goldberg \cite{popek1974}；OSTEP app.~B．}問題は，そのような
命令が特権命令であるときに何が起こらなければならないかである．

\begin{definition}[トラップアンドエミュレート]
  \term[とらっぷあんどえみゅれーと]{トラップアンドエミュレート}%
  [trap-and-emulate]による仮想化では，ゲストの命令は，ゲスト OS を降格させた
  うえで CPU 上で直接実行される．非特権命令はハードウェアの速度で動作する．
  特権命令はハイパーバイザへトラップし，ハイパーバイザが何をすべきかを判断
  する．その操作が VM にとって不正であれば VM を終了させ，正当であれば
  ゲスト OS がハードウェアに期待する振る舞いをエミュレートしてゲストを再開
  する．
\end{definition}

\cref{fig:l12-trap-emulate}はこれらの判断をまとめたものである．エミュレー
ションは，物理ハードウェアではなく VM の仮想状態に対して働く．例えばゲスト OS
が割り込みを禁止すると，ハイパーバイザはこの VM について割り込みが禁止されて
いることを記録し，その VM 宛ての割り込みを保留する．一方，物理的な割り込みは
有効なままであり，ハイパーバイザと他の VM は自分宛ての割り込みを受け取り
続ける．トラップのたびにハイパーバイザへの遷移の代価がかかるので，トラップ
アンドエミュレートのオーバーヘッドは，ゲストが行う特権的な操作の数とともに
増える．\margin{ハードウェアの遷移だけで Haswell では 500 サイクル未満，
KVM への \texttt{VMCALL} の往復は約 \num{1500} サイクルと測定されている
（kvm-unit-tests の \texttt{vmexit}）．}

\begin{figure}[htbp]
  \centering
  % Trap-and-emulate: non-privileged guest instructions run directly on the
% CPU; privileged ones trap to the hypervisor, which emulates legal
% operations and terminates the VM on illegal ones.
% Usage: % Trap-and-emulate: non-privileged guest instructions run directly on the
% CPU; privileged ones trap to the hypervisor, which emulates legal
% operations and terminates the VM on illegal ones.
% Usage: % Trap-and-emulate: non-privileged guest instructions run directly on the
% CPU; privileged ones trap to the hypervisor, which emulates legal
% operations and terminates the VM on illegal ones.
% Usage: \input{figures/l12-trap-emulate}   (width ~118 mm)
\begin{tikzpicture}[x=1mm, y=1mm, font=\scriptsize\sffamily,
  >={Stealth[length=1.5mm]},
  stp/.style={draw, rounded corners=1.5pt, fill=white, align=center,
               minimum height=9mm, inner sep=1pt},
  hyp/.style={stp, fill=ln-accent!15},
  dec/.style={draw, diamond, aspect=2, fill=ln-box, inner sep=0.3pt,
              align=center},
  yn/.style={font=\scriptsize\sffamily\itshape, text=ln-muted}]
  \node[stp, minimum width=24mm] (ins) at (12,28)
    {\iflangja{ゲストが命令を\\実行する}{guest executes\\an instruction}};
  \node[dec] (priv) at (42,28) {\iflangja{特権命令?}{privileged?}};
  \node[stp, minimum width=38mm] (hw) at (93,28)
    {\iflangja{CPU 上で直接実行\\（ハードウェアの速度）}{runs directly on the CPU\\(hardware speed)}};
  \node[hyp, minimum width=24mm] (trap) at (42,13)
    {\iflangja{ハイパーバイザへ\\トラップ}{trap to the\\hypervisor}};
  \node[dec] (legal) at (70,13) {\iflangja{正当?}{legal?}};
  \node[hyp, minimum width=30mm] (emu) at (102,13)
    {\iflangja{期待された効果を仮想\\状態上でエミュレート}{emulate the expected\\effect on virtual state}};
  \node[hyp, minimum width=26mm, minimum height=7mm] (term) at (70,-2)
    {\iflangja{VM を終了させる}{terminate the VM}};
  \draw[->] (ins) -- (priv);
  \draw[->] (priv) -- node[yn, above] {\iflangja{いいえ}{no}} (hw);
  \draw[->] (priv) -- node[yn, right] {\iflangja{はい}{yes}} (trap);
  \draw[->] (trap) -- (legal);
  \draw[->] (legal) -- node[yn, above] {\iflangja{はい}{yes}} (emu);
  \draw[->] (legal) -- node[yn, right] {\iflangja{いいえ}{no}} (term);
  \draw[->, densely dashed, ln-muted] (emu.east) -- ++(2.5,0) |- (12,38)
    -- (ins.north);
  \draw[densely dashed, ln-muted] (hw.north) -- (93,38);
  \node[yn, anchor=south] at (52,38) {\iflangja{次の命令（ゲストを再開）}{next instruction (guest resumes)}};
\end{tikzpicture}
   (width ~118 mm)
\begin{tikzpicture}[x=1mm, y=1mm, font=\scriptsize\sffamily,
  >={Stealth[length=1.5mm]},
  stp/.style={draw, rounded corners=1.5pt, fill=white, align=center,
               minimum height=9mm, inner sep=1pt},
  hyp/.style={stp, fill=ln-accent!15},
  dec/.style={draw, diamond, aspect=2, fill=ln-box, inner sep=0.3pt,
              align=center},
  yn/.style={font=\scriptsize\sffamily\itshape, text=ln-muted}]
  \node[stp, minimum width=24mm] (ins) at (12,28)
    {\iflangja{ゲストが命令を\\実行する}{guest executes\\an instruction}};
  \node[dec] (priv) at (42,28) {\iflangja{特権命令?}{privileged?}};
  \node[stp, minimum width=38mm] (hw) at (93,28)
    {\iflangja{CPU 上で直接実行\\（ハードウェアの速度）}{runs directly on the CPU\\(hardware speed)}};
  \node[hyp, minimum width=24mm] (trap) at (42,13)
    {\iflangja{ハイパーバイザへ\\トラップ}{trap to the\\hypervisor}};
  \node[dec] (legal) at (70,13) {\iflangja{正当?}{legal?}};
  \node[hyp, minimum width=30mm] (emu) at (102,13)
    {\iflangja{期待された効果を仮想\\状態上でエミュレート}{emulate the expected\\effect on virtual state}};
  \node[hyp, minimum width=26mm, minimum height=7mm] (term) at (70,-2)
    {\iflangja{VM を終了させる}{terminate the VM}};
  \draw[->] (ins) -- (priv);
  \draw[->] (priv) -- node[yn, above] {\iflangja{いいえ}{no}} (hw);
  \draw[->] (priv) -- node[yn, right] {\iflangja{はい}{yes}} (trap);
  \draw[->] (trap) -- (legal);
  \draw[->] (legal) -- node[yn, above] {\iflangja{はい}{yes}} (emu);
  \draw[->] (legal) -- node[yn, right] {\iflangja{いいえ}{no}} (term);
  \draw[->, densely dashed, ln-muted] (emu.east) -- ++(2.5,0) |- (12,38)
    -- (ins.north);
  \draw[densely dashed, ln-muted] (hw.north) -- (93,38);
  \node[yn, anchor=south] at (52,38) {\iflangja{次の命令（ゲストを再開）}{next instruction (guest resumes)}};
\end{tikzpicture}
   (width ~118 mm)
\begin{tikzpicture}[x=1mm, y=1mm, font=\scriptsize\sffamily,
  >={Stealth[length=1.5mm]},
  stp/.style={draw, rounded corners=1.5pt, fill=white, align=center,
               minimum height=9mm, inner sep=1pt},
  hyp/.style={stp, fill=ln-accent!15},
  dec/.style={draw, diamond, aspect=2, fill=ln-box, inner sep=0.3pt,
              align=center},
  yn/.style={font=\scriptsize\sffamily\itshape, text=ln-muted}]
  \node[stp, minimum width=24mm] (ins) at (12,28)
    {\iflangja{ゲストが命令を\\実行する}{guest executes\\an instruction}};
  \node[dec] (priv) at (42,28) {\iflangja{特権命令?}{privileged?}};
  \node[stp, minimum width=38mm] (hw) at (93,28)
    {\iflangja{CPU 上で直接実行\\（ハードウェアの速度）}{runs directly on the CPU\\(hardware speed)}};
  \node[hyp, minimum width=24mm] (trap) at (42,13)
    {\iflangja{ハイパーバイザへ\\トラップ}{trap to the\\hypervisor}};
  \node[dec] (legal) at (70,13) {\iflangja{正当?}{legal?}};
  \node[hyp, minimum width=30mm] (emu) at (102,13)
    {\iflangja{期待された効果を仮想\\状態上でエミュレート}{emulate the expected\\effect on virtual state}};
  \node[hyp, minimum width=26mm, minimum height=7mm] (term) at (70,-2)
    {\iflangja{VM を終了させる}{terminate the VM}};
  \draw[->] (ins) -- (priv);
  \draw[->] (priv) -- node[yn, above] {\iflangja{いいえ}{no}} (hw);
  \draw[->] (priv) -- node[yn, right] {\iflangja{はい}{yes}} (trap);
  \draw[->] (trap) -- (legal);
  \draw[->] (legal) -- node[yn, above] {\iflangja{はい}{yes}} (emu);
  \draw[->] (legal) -- node[yn, right] {\iflangja{いいえ}{no}} (term);
  \draw[->, densely dashed, ln-muted] (emu.east) -- ++(2.5,0) |- (12,38)
    -- (ins.north);
  \draw[densely dashed, ln-muted] (hw.north) -- (93,38);
  \node[yn, anchor=south] at (52,38) {\iflangja{次の命令（ゲストを再開）}{next instruction (guest resumes)}};
\end{tikzpicture}

  \caption{トラップアンドエミュレート．ハイパーバイザが関与するのは特権命令
    だけであり，正当なものは VM の仮想状態上でエミュレートされる．}
  \label{fig:l12-trap-emulate}
\end{figure}

\begin{note}
  トラップアンドエミュレートが正しく働くのは，ハイパーバイザが効果を制御し
  なければならない命令がすべて実際にトラップする場合に限られる．Popek と
  Goldberg \cite{popek1974} はそのような命令を\emph{センシティブ}命令と呼ぶ．
  マシンの構成を変更あるいは露出させる命令，または特権レベルによって振る舞いが
  変わる命令である．すべてのセンシティブ命令が特権命令である，すなわち最も
  特権的なレベルの外で実行されるとトラップするなら，そのアーキテクチャはこの
  方法で仮想化できる．
\end{note}

\section{ハードウェア支援以前の x86 の仮想化}
\label{sec:l12-x86}

2005 年以前の x86 プロセッサは4つのリングを備えていたが，ルートモードと非
ルートモードは持たなかったので，ハイパーバイザはリング 0 で動作し，ゲスト OS
をリング 1 へ降格させた．\source{P3L6，過去の x86 の仮想化；Robin と Irvine
\cite{robin2000}．}このアーキテクチャはトラップアンドエミュレートの条件を
満たしていなかった．Robin と Irvine \cite{robin2000} は，特権の低いリングで
実行してもトラップしないセンシティブ命令を 17 個特定した．\caution{この手は
64 ビットのゲストには通じない．AMD64 のロングモードはセグメント限界の検査を
やめたので，ソフトウェアの VMM はセグメンテーションでゲストを囲い込めない．
\cref{sec:l12-hw}の支援が要る．}リング 1 で実行する
と，それらは黙って失敗するか，ハイパーバイザが隠そうとしているマシンの状態を
露出させるので，トラップアンドエミュレートだけではプロセッサを正しく仮想化
できない．

最もよく知られた例は \texttt{PUSHF} と \texttt{POPF} であり，フラグレジスタを
スタックに保存し，また復元する．フラグレジスタには割り込み許可フラグ IF が
含まれており，カーネルは，割り込みを禁止して実行した区間の終わりで，保存した
フラグを \texttt{POPF} で復元することによって割り込みを再び許可するのが普通で
ある．リング 0 では \texttt{POPF} は IF を復元する．リング 1 では，プロセッサ
は復元する値の IF ビットを，トラップすることなく無視する．\margin{ハイパー
バイザは I/O 特権レベルを 0 に設定するので，リング 1 の \texttt{POPF} は IF を
変更しないままにし，一方 \texttt{CLI} はそこでトラップする．}ハイパーバイザは
呼び出されないので，ゲストが割り込みの状態を変えようとしたことを知ることが
できず，その変更をエミュレートできない．ゲスト OS はエラーを受け取らないので，
変更は成功したものと思い込む．\cref{tab:l12-popf}はその帰結を追ったもので
ある．

\begin{table}[htbp]
  \centering
  \caption{リング 1 のゲストカーネルが，フラグを保存し，割り込みを禁止し，
    後でフラグを復元する．\texttt{CLI} はトラップしてエミュレートされるが，
    \texttt{POPF} はトラップしない．\emph{仮想 IF} は，ハイパーバイザがこの
    VM について記録している割り込みフラグであり，物理的な割り込みは終始有効な
    ままである．}
  \label{tab:l12-popf}
  \small
  \begin{tabular}{@{}clccc@{}}
    \toprule
    ステップ & ゲストカーネルの実行 & トラップ & 仮想 IF & ゲストの想定 \\
    \midrule
    1 & \texttt{PUSHF}（IF $=1$ を保存）   & しない & 1 & 許可 \\
    2 & \texttt{CLI}                       & する   & 0 & 禁止 \\
    3 & クリティカルセクション             & しない & 0 & 禁止 \\
    4 & \texttt{POPF}（IF $=1$ を復元）     & しない & 0 & 許可 \\
    5 & デバイス割り込みを待つ             & しない & 0 & 許可 \\
    \bottomrule
  \end{tabular}
\end{table}

ステップ 4 以降，ハイパーバイザはこの VM について割り込みが禁止されていると
考えてその割り込みを保留し，一方ゲスト OS は，許可されていると信じている割り
込みを待ち続ける．ゲストはハングする．これに対する対策がソフトウェアで2つ
開発された．1つはゲスト OS を改変しないもの，もう1つは改変するものである．

\section{バイナリ変換}
\label{sec:l12-bt}

第1の対策は，ゲストのバイナリコードを書き換えて，問題のある命令を決して実行
しないようにする．\source{P3L6，バイナリ変換；Bugnion ら \cite{bugnion2012}；
Silberschatz ら ch.~18．}この手法はスタンフォード大学の Mendel Rosenblum の
グループが先鞭をつけ，VMware が製品化した \cite{bugnion2012}．その目標は，
ゲスト OS をまったく改変せずに動かすことである．

\begin{definition}[完全仮想化]
  ゲストオペレーティングシステムが，物理マシン上で動作するのと同じように，
  改変されずに動作する仮想化を
  \term[かんぜんかそうか]{完全仮想化}[full virtualization]と呼ぶ．
\end{definition}

\begin{definition}[バイナリ変換]
  動的\term[ばいなりへんかん]{バイナリ変換}[binary translation]では，
  ハイパーバイザが実行時に，実行の直前にゲストのコードを書き換える．問題の
  ある命令は，意図した効果を持つ別の命令列に変換され，それ以外の命令はその
  まま残されてハードウェアの速度で動作する．
\end{definition}

変換はブロック単位で進む（\cref{fig:l12-bt}）．
\begin{enumerate}
  \item これから実行されるコードのブロック，例えば次の分岐までの命令列を検査
    する．
  \item ブロックに問題のある命令が含まれていれば，望ましい振る舞いをエミュ
    レートする命令列，例えばハイパーバイザのエミュレーションルーチンの呼び
    出しに変換する．元の命令がトラップしたはずの箇所でも，変換によってトラップ
    を避けられることさえある．
  \item そうでなければ，制御の移動だけを調整してブロックをそのまま引き取る．
    こうしてブロックはハードウェアの速度で動作しつつ，実行は変換器の制御下に
    留まる．
\end{enumerate}
得られたブロックは\emph{変換キャッシュ}に保持されるので，ブロックを変換する
代価は一度だけ支払われ，その後の実行にわたって償却される．どのコードが実行
されるか，間接分岐がどこへ向かうかは実行時にしか分からないので，変換は動的で
なければならない．%
\margin{VMware の設計では，変換されるのはゲストのカーネルコードだけであり，
リング 3 のアプリケーションは直接動作する \cite{bugnion2012}．}

\begin{figure}[htbp]
  \centering
  % Dynamic binary translation with a translation cache: a guest kernel code
% block is inspected and translated the first time it is about to run;
% later executions reuse the cached translation.
% Usage: % Dynamic binary translation with a translation cache: a guest kernel code
% block is inspected and translated the first time it is about to run;
% later executions reuse the cached translation.
% Usage: % Dynamic binary translation with a translation cache: a guest kernel code
% block is inspected and translated the first time it is about to run;
% later executions reuse the cached translation.
% Usage: \input{figures/l12-binary-translation}   (width ~119 mm)
\begin{tikzpicture}[x=1mm, y=1mm, font=\scriptsize\sffamily,
  >={Stealth[length=1.5mm]},
  stp/.style={draw, rounded corners=1.5pt, fill=white, align=center,
               minimum height=9mm, inner sep=1pt},
  hyp/.style={stp, fill=ln-accent!15},
  dec/.style={draw, diamond, aspect=2.2, fill=ln-box, inner sep=0.3pt,
              align=center},
  yn/.style={font=\scriptsize\sffamily\itshape, text=ln-muted}]
  \node[stp, minimum width=22mm] (blk) at (11,26)
    {\iflangja{次のゲスト\\カーネルコード\\ブロック}{next guest\\kernel code\\block}};
  \node[dec] (hit) at (42,26)
    {\iflangja{変換キャッシュ\\にある?}{in translation\\cache?}};
  \node[stp, minimum width=32mm] (exe) at (103,26)
    {\iflangja{変換済みブロックを\\ハードウェアの速度で実行}{execute translated block\\at hardware speed}};
  \node[hyp, minimum width=22mm] (insp) at (42,8)
    {\iflangja{命令を\\検査する}{inspect the\\instructions}};
  \node[hyp, minimum width=26mm] (tr) at (72,8)
    {\iflangja{問題のある命令を\\別の命令列に置換}{replace problematic\\instructions}};
  \node[hyp, minimum width=28mm] (st) at (103,8)
    {\iflangja{変換キャッシュに\\格納する}{store the block in\\the translation cache}};
  \draw[->] (blk) -- (hit);
  \draw[->] (hit) -- node[yn, above] {\iflangja{はい}{yes}} (exe);
  \draw[->] (hit) -- node[yn, right] {\iflangja{いいえ}{no}} (insp);
  \draw[->] (insp) -- (tr);
  \draw[->] (tr) -- (st);
  \draw[->] (st) -- (exe);
  \draw[->, densely dashed, ln-muted] (exe.north) -- (103,39) -- (11,39)
    -- (blk.north);
  \node[yn, anchor=south] at (57,39) {\iflangja{次のブロックへ}{continue with the next block}};
\end{tikzpicture}
   (width ~119 mm)
\begin{tikzpicture}[x=1mm, y=1mm, font=\scriptsize\sffamily,
  >={Stealth[length=1.5mm]},
  stp/.style={draw, rounded corners=1.5pt, fill=white, align=center,
               minimum height=9mm, inner sep=1pt},
  hyp/.style={stp, fill=ln-accent!15},
  dec/.style={draw, diamond, aspect=2.2, fill=ln-box, inner sep=0.3pt,
              align=center},
  yn/.style={font=\scriptsize\sffamily\itshape, text=ln-muted}]
  \node[stp, minimum width=22mm] (blk) at (11,26)
    {\iflangja{次のゲスト\\カーネルコード\\ブロック}{next guest\\kernel code\\block}};
  \node[dec] (hit) at (42,26)
    {\iflangja{変換キャッシュ\\にある?}{in translation\\cache?}};
  \node[stp, minimum width=32mm] (exe) at (103,26)
    {\iflangja{変換済みブロックを\\ハードウェアの速度で実行}{execute translated block\\at hardware speed}};
  \node[hyp, minimum width=22mm] (insp) at (42,8)
    {\iflangja{命令を\\検査する}{inspect the\\instructions}};
  \node[hyp, minimum width=26mm] (tr) at (72,8)
    {\iflangja{問題のある命令を\\別の命令列に置換}{replace problematic\\instructions}};
  \node[hyp, minimum width=28mm] (st) at (103,8)
    {\iflangja{変換キャッシュに\\格納する}{store the block in\\the translation cache}};
  \draw[->] (blk) -- (hit);
  \draw[->] (hit) -- node[yn, above] {\iflangja{はい}{yes}} (exe);
  \draw[->] (hit) -- node[yn, right] {\iflangja{いいえ}{no}} (insp);
  \draw[->] (insp) -- (tr);
  \draw[->] (tr) -- (st);
  \draw[->] (st) -- (exe);
  \draw[->, densely dashed, ln-muted] (exe.north) -- (103,39) -- (11,39)
    -- (blk.north);
  \node[yn, anchor=south] at (57,39) {\iflangja{次のブロックへ}{continue with the next block}};
\end{tikzpicture}
   (width ~119 mm)
\begin{tikzpicture}[x=1mm, y=1mm, font=\scriptsize\sffamily,
  >={Stealth[length=1.5mm]},
  stp/.style={draw, rounded corners=1.5pt, fill=white, align=center,
               minimum height=9mm, inner sep=1pt},
  hyp/.style={stp, fill=ln-accent!15},
  dec/.style={draw, diamond, aspect=2.2, fill=ln-box, inner sep=0.3pt,
              align=center},
  yn/.style={font=\scriptsize\sffamily\itshape, text=ln-muted}]
  \node[stp, minimum width=22mm] (blk) at (11,26)
    {\iflangja{次のゲスト\\カーネルコード\\ブロック}{next guest\\kernel code\\block}};
  \node[dec] (hit) at (42,26)
    {\iflangja{変換キャッシュ\\にある?}{in translation\\cache?}};
  \node[stp, minimum width=32mm] (exe) at (103,26)
    {\iflangja{変換済みブロックを\\ハードウェアの速度で実行}{execute translated block\\at hardware speed}};
  \node[hyp, minimum width=22mm] (insp) at (42,8)
    {\iflangja{命令を\\検査する}{inspect the\\instructions}};
  \node[hyp, minimum width=26mm] (tr) at (72,8)
    {\iflangja{問題のある命令を\\別の命令列に置換}{replace problematic\\instructions}};
  \node[hyp, minimum width=28mm] (st) at (103,8)
    {\iflangja{変換キャッシュに\\格納する}{store the block in\\the translation cache}};
  \draw[->] (blk) -- (hit);
  \draw[->] (hit) -- node[yn, above] {\iflangja{はい}{yes}} (exe);
  \draw[->] (hit) -- node[yn, right] {\iflangja{いいえ}{no}} (insp);
  \draw[->] (insp) -- (tr);
  \draw[->] (tr) -- (st);
  \draw[->] (st) -- (exe);
  \draw[->, densely dashed, ln-muted] (exe.north) -- (103,39) -- (11,39)
    -- (blk.north);
  \node[yn, anchor=south] at (57,39) {\iflangja{次のブロックへ}{continue with the next block}};
\end{tikzpicture}

  \caption{変換キャッシュを伴う動的バイナリ変換．ブロックが検査され変換され
    るのは，最初に実行されようとするときだけである．}
  \label{fig:l12-bt}
\end{figure}

\needspace{8\baselineskip}
\begin{example}
  ゲストのカーネルコードのあるブロックは，ネイティブには
  \SI{5}{\micro\second} で実行される．\sidenote{バイナリ変換は，ハードウェア
  支援が現れても時代遅れにはならなかった．最初の VT-x プロセッサでは，VM exit
  が高価で，どちらにせよページテーブルをシャドウ化しなければならず，ほとんどの
  ワークロードではソフトウェアの VMM の方が速かった \cite{adams2006}．}%
  その変換には \SI{2}{\milli\second}
  かかり，1 命令がハイパーバイザへの呼び出しに置き換えられるため，変換後の
  ブロックの実行には \SI{6}{\micro\second} かかる．このブロックが
  \num{20000} 回実行されるなら，一度だけ変換して結果をキャッシュする代価は
  $\SI{2}{\milli\second} + \num{20000}\times\SI{6}{\micro\second} =
  \SI{122}{\milli\second}$ であり，ネイティブの \SI{100}{\milli\second} に
  対してオーバーヘッドは \SI{22}{\percent} である．実行のたびに変換し直せば
  $\num{20000}\times(\SI{2}{\milli\second}+\SI{6}{\micro\second}) \approx
  \SI{40.1}{\second}$ となり，ネイティブの約 400 倍の時間になる．
\end{example}

\section{準仮想化}
\label{sec:l12-pv}

第2の対策は，性能と引き換えに，改変しないゲストを諦める．\source{P3L6，
準仮想化；Barham ら \cite{barham2003}．}ゲスト OS は，オペレーティング
システムを新しいハードウェアプラットフォームへ移植するのとほぼ同じように，
ハイパーバイザへ移植される．

\begin{definition}[準仮想化]
  \term[じゅんかそうか]{準仮想化}[paravirtualization]では，ゲストオペレー
  ティングシステムが，自分は VM の中で動作していると知るように改変される．
  特権的な操作を実行してトラップに頼るのではなく，それをハイパーバイザに明示
  的に要求する．
\end{definition}

ゲスト OS からハイパーバイザへの明示的な要求が
\term[はいぱーこーる]{ハイパーコール}[hypercall]である．これは，アプリ
ケーションにとってのシステムコール（第1章）と同じ役割を，ゲスト OS に対して
果たす．ゲストは文脈情報---引数---を定められた場所に置き，どのハイパーコール
を要求するかを指定してハイパーバイザへトラップする．ハイパーバイザはその操作
を実行し，制御をゲストへ返す．したがってハイパーバイザは，トラップした個々の
命令から意図された操作を推測する必要がなく，\cref{sec:l12-memory}で見る
ように，1回のハイパーコールで複数の操作を運べる．改変されるのはゲスト OS
だけであり，ゲスト OS がアプリケーションに提供するインタフェースは変わらない
ので，アプリケーションはそのまま動作する \cite{barham2003}．

Xen はケンブリッジ大学で開発され，後に XenSource と Citrix によって製品化
された，最もよく知られた準仮想化ハイパーバイザである．\cref{sec:l12-hw}の
\cref{tab:l12-cpu}は，準仮想化とプロセッサ仮想化の他の手法を比較したもので
ある．

\begin{keypoint}
  トラップアンドエミュレートは，すべてのセンシティブ命令がトラップすることを
  必要とするが，2005 年以前の x86 プロセッサはそれを保証しなかった．バイナリ
  変換は問題のある命令を実行時に書き換え，ゲスト OS を改変しないまま保つ．
  準仮想化はゲスト OS を改変してハイパーコールを発行させ，性能を得る．
\end{keypoint}

\section{メモリの仮想化}
\label{sec:l12-memory}

\subsection{完全仮想化}

ゲスト OS は，オペレーティングシステムが物理マシン上で見出すもの，すなわち
アドレス 0 から始まる連続した物理メモリを期待し，第8章で述べたページテーブル
でそれを管理する．\source{P3L6，メモリの仮想化: 完全仮想化，準仮想化；
OSTEP app.~B；Barham ら \cite{barham2003}．}ハイパーバイザは，マシンのメモリ
を VM の間で分割しながら，すべての VM にこの幻想を与えなければならず，しかも
アドレス変換はハードウェアの MMU と TLB に行わせたい．そこで3種類のアドレスを
区別する．
\begin{description}
  \item[仮想アドレス] 第2章と同じく，VM 内のソフトウェアが使うアドレスで
    ある．
  \item[物理アドレス] ゲスト OS が物理メモリのアドレスだと信じているアドレス
    である．
  \item[マシンアドレス] マシンのメモリの実際のアドレスである．完全仮想化では，
    \term[ましんあどれす]{マシンアドレス}[machine address]はハイパーバイザ
    だけが知っている．
\end{description}
ページ番号についても同じ区別があり，仮想ページ番号，物理フレーム番号，マシン
フレーム番号となる．ゲスト OS は仮想アドレスを物理アドレスへ対応付けるページ
テーブルを維持し，ハイパーバイザは物理アドレスからマシンアドレスへの対応を
維持する（\cref{fig:l12-memory}）．\tip{2つの矢印を順に読めばよい．仮想
$\to$ 物理はゲストのもの，物理 $\to$ マシンはハイパーバイザのものである．
シャドウページテーブルはその合成，すなわち仮想 $\to$ マシンである．}

\begin{figure}[htbp]
  \centering
  % Memory virtualization: virtual, physical and machine addresses.  The
% guest page table maps VA to PA, the hypervisor maps PA to MA; the shadow
% page table maps VA directly to MA and is the table the MMU uses.
% Usage: % Memory virtualization: virtual, physical and machine addresses.  The
% guest page table maps VA to PA, the hypervisor maps PA to MA; the shadow
% page table maps VA directly to MA and is the table the MMU uses.
% Usage: % Memory virtualization: virtual, physical and machine addresses.  The
% guest page table maps VA to PA, the hypervisor maps PA to MA; the shadow
% page table maps VA directly to MA and is the table the MMU uses.
% Usage: \input{figures/l12-memory}   (width ~120 mm)
\begin{tikzpicture}[x=1mm, y=1mm, font=\scriptsize\sffamily,
  >={Stealth[length=1.5mm]},
  addr/.style={draw, rounded corners=4pt, fill=white, align=center,
               minimum height=9mm, inner sep=1pt},
  tab/.style={draw, fill=ln-accent!12, align=center, minimum height=10mm,
              minimum width=22mm, inner sep=1pt},
  htab/.style={tab, fill=ln-accent!30},
  own/.style={font=\scriptsize\sffamily\bfseries, text=ln-accent},
  lab/.style={font=\scriptsize\sffamily, text=ln-muted, align=center}]
  \node[addr, minimum width=17mm] (va) at (8.5,22)
    {\iflangja{仮想\\アドレス}{virtual\\address}};
  \node[tab] (gpt) at (34,22)
    {\iflangja{ゲストの\\ページテーブル}{guest\\page table}};
  \node[addr, minimum width=20mm] (pa) at (60,22)
    {\iflangja{物理\\アドレス}{physical\\address}};
  \node[htab] (p2m) at (86,22)
    {\iflangja{ハイパーバイザの\\PA$\to$MA 対応表}{hypervisor's\\PA$\to$MA map}};
  \node[addr, minimum width=17mm] (ma) at (111.5,22)
    {\iflangja{マシン\\アドレス}{machine\\address}};
  \draw[->] (va) -- (gpt);
  \draw[->] (gpt) -- (pa);
  \draw[->] (pa) -- (p2m);
  \draw[->] (p2m) -- (ma);
  \node[own, anchor=south] at (gpt.north) {\iflangja{ゲスト OS が管理}{kept by the guest OS}};
  \node[own, anchor=south] at (p2m.north) {\iflangja{ハイパーバイザが管理}{kept by the hypervisor}};
  \draw[ln-muted] (23,31.5) -- (23,33.5) -- (97,33.5) -- (97,31.5);
  \node[lab, anchor=south] at (60,33.7) {\iflangja{2 回の変換}{two translations per address}};
  % shadow page table
  \node[htab, minimum width=40mm, minimum height=8mm] (spt) at (60,4)
    {\iflangja{シャドウページテーブル（VA$\to$MA）}{shadow page table (VA$\to$MA)}};
  \draw[->, very thick, ln-accent] (va.south) |- (spt.west);
  \draw[->, very thick, ln-accent] (spt.east) -| (ma.south);
  \draw[->, densely dashed, ln-muted] (gpt.south) -- (44,8.3);
  \draw[->, densely dashed, ln-muted] (p2m.south) -- (76,8.3);
  \node[lab] at (60,12.8) {\iflangja{ハイパーバイザが一貫性を保つ}{kept consistent by the hypervisor}};
  \node[lab, anchor=north] at (60,-0.4)
    {\iflangja{1 回の変換．MMU と TLB が使用}{one translation, used by the MMU and TLB}};
\end{tikzpicture}
   (width ~120 mm)
\begin{tikzpicture}[x=1mm, y=1mm, font=\scriptsize\sffamily,
  >={Stealth[length=1.5mm]},
  addr/.style={draw, rounded corners=4pt, fill=white, align=center,
               minimum height=9mm, inner sep=1pt},
  tab/.style={draw, fill=ln-accent!12, align=center, minimum height=10mm,
              minimum width=22mm, inner sep=1pt},
  htab/.style={tab, fill=ln-accent!30},
  own/.style={font=\scriptsize\sffamily\bfseries, text=ln-accent},
  lab/.style={font=\scriptsize\sffamily, text=ln-muted, align=center}]
  \node[addr, minimum width=17mm] (va) at (8.5,22)
    {\iflangja{仮想\\アドレス}{virtual\\address}};
  \node[tab] (gpt) at (34,22)
    {\iflangja{ゲストの\\ページテーブル}{guest\\page table}};
  \node[addr, minimum width=20mm] (pa) at (60,22)
    {\iflangja{物理\\アドレス}{physical\\address}};
  \node[htab] (p2m) at (86,22)
    {\iflangja{ハイパーバイザの\\PA$\to$MA 対応表}{hypervisor's\\PA$\to$MA map}};
  \node[addr, minimum width=17mm] (ma) at (111.5,22)
    {\iflangja{マシン\\アドレス}{machine\\address}};
  \draw[->] (va) -- (gpt);
  \draw[->] (gpt) -- (pa);
  \draw[->] (pa) -- (p2m);
  \draw[->] (p2m) -- (ma);
  \node[own, anchor=south] at (gpt.north) {\iflangja{ゲスト OS が管理}{kept by the guest OS}};
  \node[own, anchor=south] at (p2m.north) {\iflangja{ハイパーバイザが管理}{kept by the hypervisor}};
  \draw[ln-muted] (23,31.5) -- (23,33.5) -- (97,33.5) -- (97,31.5);
  \node[lab, anchor=south] at (60,33.7) {\iflangja{2 回の変換}{two translations per address}};
  % shadow page table
  \node[htab, minimum width=40mm, minimum height=8mm] (spt) at (60,4)
    {\iflangja{シャドウページテーブル（VA$\to$MA）}{shadow page table (VA$\to$MA)}};
  \draw[->, very thick, ln-accent] (va.south) |- (spt.west);
  \draw[->, very thick, ln-accent] (spt.east) -| (ma.south);
  \draw[->, densely dashed, ln-muted] (gpt.south) -- (44,8.3);
  \draw[->, densely dashed, ln-muted] (p2m.south) -- (76,8.3);
  \node[lab] at (60,12.8) {\iflangja{ハイパーバイザが一貫性を保つ}{kept consistent by the hypervisor}};
  \node[lab, anchor=north] at (60,-0.4)
    {\iflangja{1 回の変換．MMU と TLB が使用}{one translation, used by the MMU and TLB}};
\end{tikzpicture}
   (width ~120 mm)
\begin{tikzpicture}[x=1mm, y=1mm, font=\scriptsize\sffamily,
  >={Stealth[length=1.5mm]},
  addr/.style={draw, rounded corners=4pt, fill=white, align=center,
               minimum height=9mm, inner sep=1pt},
  tab/.style={draw, fill=ln-accent!12, align=center, minimum height=10mm,
              minimum width=22mm, inner sep=1pt},
  htab/.style={tab, fill=ln-accent!30},
  own/.style={font=\scriptsize\sffamily\bfseries, text=ln-accent},
  lab/.style={font=\scriptsize\sffamily, text=ln-muted, align=center}]
  \node[addr, minimum width=17mm] (va) at (8.5,22)
    {\iflangja{仮想\\アドレス}{virtual\\address}};
  \node[tab] (gpt) at (34,22)
    {\iflangja{ゲストの\\ページテーブル}{guest\\page table}};
  \node[addr, minimum width=20mm] (pa) at (60,22)
    {\iflangja{物理\\アドレス}{physical\\address}};
  \node[htab] (p2m) at (86,22)
    {\iflangja{ハイパーバイザの\\PA$\to$MA 対応表}{hypervisor's\\PA$\to$MA map}};
  \node[addr, minimum width=17mm] (ma) at (111.5,22)
    {\iflangja{マシン\\アドレス}{machine\\address}};
  \draw[->] (va) -- (gpt);
  \draw[->] (gpt) -- (pa);
  \draw[->] (pa) -- (p2m);
  \draw[->] (p2m) -- (ma);
  \node[own, anchor=south] at (gpt.north) {\iflangja{ゲスト OS が管理}{kept by the guest OS}};
  \node[own, anchor=south] at (p2m.north) {\iflangja{ハイパーバイザが管理}{kept by the hypervisor}};
  \draw[ln-muted] (23,31.5) -- (23,33.5) -- (97,33.5) -- (97,31.5);
  \node[lab, anchor=south] at (60,33.7) {\iflangja{2 回の変換}{two translations per address}};
  % shadow page table
  \node[htab, minimum width=40mm, minimum height=8mm] (spt) at (60,4)
    {\iflangja{シャドウページテーブル（VA$\to$MA）}{shadow page table (VA$\to$MA)}};
  \draw[->, very thick, ln-accent] (va.south) |- (spt.west);
  \draw[->, very thick, ln-accent] (spt.east) -| (ma.south);
  \draw[->, densely dashed, ln-muted] (gpt.south) -- (44,8.3);
  \draw[->, densely dashed, ln-muted] (p2m.south) -- (76,8.3);
  \node[lab] at (60,12.8) {\iflangja{ハイパーバイザが一貫性を保つ}{kept consistent by the hypervisor}};
  \node[lab, anchor=north] at (60,-0.4)
    {\iflangja{1 回の変換．MMU と TLB が使用}{one translation, used by the MMU and TLB}};
\end{tikzpicture}

  \caption{メモリの仮想化．ゲストのページテーブルを通り，次にハイパーバイザの
    対応表を通る変換は2段階を要する．シャドウページテーブルは両者を，MMU が
    直接使う1つの対応にまとめる．拡張ページテーブルを備えたハードウェアは，
    この2段階の変換を自ら行う（\cref{sec:l12-hw}）．}
  \label{fig:l12-memory}
\end{figure}

第1の選択肢は，すべてのアドレスをこの2段階で変換することである．しかし MMU が
辿るページテーブルは1つだけであり，TLB はアドレスごとに1つの変換しかキャッシュ
しない．メモリ参照のたびに第2の変換をソフトウェアで行うのは，あまりにも高価で
ある．

\begin{definition}[シャドウページテーブル]
  \term[しゃどうぺーじてーぶる]{シャドウページテーブル}[shadow page table]と
  は，ハイパーバイザが維持するページテーブルで，ゲストプロセスの仮想アドレスを
  マシンアドレスへ直接対応付ける．ハイパーバイザはこれをゲストのページテーブル
  の代わりに MMU に設定するので，ハードウェアは各アドレスを1段階で変換し，TLB
  がその結果をキャッシュする．
\end{definition}

ゲスト OS は自分のページテーブルを作り，変更し続けるが，ハードウェアはもはや
それを使わない．そのためハイパーバイザは，シャドウページテーブルをゲストの
ページテーブルと整合させ続けなければならない．
\begin{itemize}
  \item ゲストのページテーブルを保持するメモリを書き込み禁止にする．ゲスト OS
    が対応付けを追加・変更すると書き込みがトラップし，ハイパーバイザがその
    書き込みを行ったうえで，対応するマシンアドレスでシャドウページテーブルを
    更新する．\margin{書き込みを知るためにページを書き込み禁止にする仕掛けは，
    第8章のコピーオンライトと同じである．ここではページを複製するためでは
    なく，ゲストを観察するために使う．}
  \item ゲスト OS がプロセスを切り替えるときは，別のページテーブルをページ
    テーブルベースレジスタに読み込む．これは特権的な操作なのでトラップする．
    そこでハイパーバイザは，新しいプロセスのシャドウページテーブルを設定する
    か，現在のシャドウエントリを無効化して，使われるにつれて作り直す．
\end{itemize}
これらのトラップによって，ゲスト内でのページテーブルの更新とコンテキスト
スイッチは，物理マシン上よりはるかに高価になる．

\begin{example}
  あるゲストプロセスが仮想ページ 2 と 3 を使っている．ゲストのページテーブルは
  それらを物理フレーム 5 と 9 に対応付けており，ハイパーバイザはそれらを
  マシンフレーム 41 と 17 で裏付けているので，シャドウページテーブルはページ 2
  をマシンフレーム 41 に，ページ 3 をマシンフレーム 17 に対応付ける．ここで
  ゲスト OS が仮想ページ 4 を物理フレーム 12 に対応付ける．保護されたページ
  テーブルへの書き込みはトラップし，ハイパーバイザはフレーム 12 がマシン
  フレーム 26 で裏付けられていることを見つけ，書き込みを行い，シャドウページ
  テーブルにエントリ $4 \to 26$ を追加してゲストを再開する．
  \cref{tab:l12-shadow}に，その結果の対応を示す．
\end{example}

\begin{table}[htbp]
  \centering
  \caption{例のゲストプロセスの対応．最後のシャドウエントリは，ゲストの
    ページテーブル更新がトラップしたときにハイパーバイザが追加したもので
    ある．}
  \label{tab:l12-shadow}
  \small
  \begin{tabular}{@{}cccc@{}}
    \toprule
    仮想 & ゲストのページテーブル: & ハイパーバイザ: & シャドウページテーブル: \\
    ページ & 物理フレーム & マシンフレーム & マシンフレーム \\
    \midrule
    2 & 5  & 41 & 41 \\
    3 & 9  & 17 & 17 \\
    4 & 12 & 26 & 26 \\
    \bottomrule
  \end{tabular}
\end{table}

\subsection{準仮想化}

準仮想化ではゲスト OS は自分が VM の中で動作していることを知っており，メモリの
仮想化はゲストとハイパーバイザの協調になる．物理メモリが連続していてアドレス 0
から始まるという厳しい要求は消える．Xen では，ゲスト OS は自分に与えられた
マシンフレームを知っており，そのページテーブルは仮想アドレスをマシンアドレスへ
直接対応付ける \cite{barham2003}．\caution{呼び方はシステムごとに異なる．
この章と Xen が物理アドレス，マシンアドレスと呼ぶものを，Intel はゲスト物理
アドレス，ホスト物理アドレスと呼び，VMware も同様に呼ぶ．}%
ゲストは自分のページテーブルをハイパーバイザ
に明示的に登録し，それ以降そのページテーブルに対しては読み出しの権限しか
持たない．ページテーブルを変更するにはハイパーコールを発行し，ハイパーバイザは
その変更を---例えば新しい対応がこの VM に属するマシンフレームを指していること
を---検証してから適用する．ハイパーコールのたびにハイパーバイザへの遷移の代価
がかかるので，ゲストはページテーブルの更新を\emph{まとめ}，多数の更新が1回の
ハイパーコールを共有するようにする．他の最適化も，遷移を減らすという同じ考え
に従う．例えば Xen は，すべてのゲストのアドレス空間の上位
\SI{64}{\mega\byte} を占めるので，ハイパーバイザに入るのにページテーブルの
切り替えも TLB のフラッシュも必要ない \cite{barham2003}．

\needspace{16\baselineskip}
\begin{example}
  Xen 上の準仮想化された Linux ゲストが，1回のハイパーコールで3つのページ
  テーブルエントリを変更する．各要求は，ページテーブルエントリのマシンアドレス
  と，そのエントリの新しい内容を与える:
  \begin{lstlisting}[language=C, numbers=none]
struct mmu_update req[3];
int done;

for (int i = 0; i < 3; i++) {
  /* 変更するエントリのマシンアドレス */
  req[i].ptr = pte_maddr[i] | MMU_NORMAL_PT_UPDATE;
  /* 新しいエントリ: マシンフレーム番号とフラグ */
  req[i].val = (mfn[i] << PAGE_SHIFT) | flags;
}
/* 3 つの更新すべてに対して Xen への遷移は 1 回 */
HYPERVISOR_mmu_update(req, 3, &done, DOMID_SELF);
\end{lstlisting}
  Xen は各要求を適用する前に検査し，いくつの更新が成功したかを
  \texttt{done} に報告する．1つずつ発行すれば，同じ更新のためにハイパーバイザ
  へ3回入ることになる．
\end{example}

より新しいプラットフォームでは，\cref{sec:l12-hw}のハードウェア支援が，メモリ
の仮想化のオーバーヘッドを取り除くか大幅に減らす．

\section{デバイスの仮想化}
\label{sec:l12-devices}

CPU とメモリの多様性は小さい．命令セットアーキテクチャがそのインタフェースを
標準化しているので，ある系列のすべてのプロセッサにわずかな技法で足りる．
\source{P3L6，デバイスの仮想化，パススルー，ハイパーバイザ直接，分割デバイス
ドライバの各モデル；Silberschatz ら ch.~18．}デバイスは非常に多様であり，その
インタフェースと振る舞いの標準仕様は存在しない．第11章で見たように，デバイスの
種類ごとに専用のドライバが必要である．デバイスを仮想化するモデルは3つ現れた．
物理デバイスを制御するドライバがどこに置かれるか，VM のデバイスアクセスがどの
ようにそこへ届くかが異なる（\cref{fig:l12-devices}）．

\begin{figure}[htbp]
  \centering
  % Device virtualization models: (a) passthrough, (b) hypervisor-direct,
% (c) split-device driver with front-end and back-end drivers.
% Usage: % Device virtualization models: (a) passthrough, (b) hypervisor-direct,
% (c) split-device driver with front-end and back-end drivers.
% Usage: % Device virtualization models: (a) passthrough, (b) hypervisor-direct,
% (c) split-device driver with front-end and back-end drivers.
% Usage: \input{figures/l12-devices}   (width ~118 mm)
\begin{tikzpicture}[x=1mm, y=1mm, font=\scriptsize\sffamily,
  >={Stealth[length=1.5mm]},
  dev/.style={draw, fill=white},
  band/.style={draw, fill=ln-accent!22},
  vm/.style={draw=ln-rule, rounded corners=2pt, fill=ln-box},
  svm/.style={draw=ln-accent, rounded corners=2pt, fill=ln-accent!8},
  gos/.style={draw, fill=ln-accent!12},
  box/.style={draw, fill=white, align=center, inner sep=0.5pt,
              font=\scriptsize\sffamily},
  hdr/.style={font=\scriptsize\sffamily\bfseries},
  lab/.style={font=\scriptsize\sffamily, text=ln-muted},
  io/.style={->, thick, ln-accent},
  pcap/.style={font=\footnotesize\sffamily, anchor=north}]
  % ================= (a) passthrough
  \draw[dev] (0,0) rectangle (36,6) node[midway] {\iflangja{デバイス}{device}};
  \draw[band] (0,8.5) rectangle (22,29);
  \node[hdr] at (11,26.5) {\iflangja{ハイパーバイザ}{hypervisor}};
  \node[box, minimum width=18mm, minimum height=8mm] (vd) at (11,16.5)
    {\iflangja{VMM レベル\\ドライバ}{VMM-level\\driver}};
  \draw[->, densely dashed] (vd.south) -- (11,6);
  \draw[vm] (0,31) rectangle (36,55);
  \node[hdr] at (18,52.5) {\iflangja{ゲスト VM}{guest VM}};
  \draw[gos] (2,33) rectangle (34,49.5);
  \node[lab, text=black] at (18,47.2) {\iflangja{ゲスト OS}{guest OS}};
  \node[box, minimum width=24mm, minimum height=7mm] at (18,39.5)
    {\iflangja{デバイスドライバ}{device driver}};
  \draw[io] (29,36) -- (29,6);
  \node[lab, rotate=90, text=ln-accent] at (32.2,19) {\iflangja{直接アクセス}{direct access}};
  \node[pcap] at (18,-2) {\iflangja{(a) パススルー}{(a) passthrough}};
  % ================= (b) hypervisor-direct
  \begin{scope}[xshift=41mm]
  \draw[dev] (0,0) rectangle (36,6) node[midway] {\iflangja{デバイス}{device}};
  \draw[band] (0,8.5) rectangle (36,29);
  \node[hdr, anchor=west] at (1.2,26.5) {\iflangja{ハイパーバイザ}{hypervisor}};
  \node[box, minimum width=30mm, minimum height=4.3mm] (em) at (18,21.7)
    {\iflangja{デバイスエミュレーション}{device emulation}};
  \node[box, minimum width=30mm, minimum height=4.3mm] (st) at (18,16.6)
    {\iflangja{I/O スタック}{I/O stack}};
  \node[box, minimum width=30mm, minimum height=4.3mm] (dd) at (18,11.5)
    {\iflangja{デバイスドライバ}{device driver}};
  \draw[vm] (0,31) rectangle (36,55);
  \node[hdr] at (18,52.5) {\iflangja{ゲスト VM}{guest VM}};
  \draw[gos] (2,33) rectangle (34,49.5);
  \node[lab, text=black] at (18,47.2) {\iflangja{ゲスト OS}{guest OS}};
  \node[box, minimum width=26mm, minimum height=7mm] (gd) at (18,39.5)
    {\iflangja{エミュレートされた\\デバイスのドライバ}{driver for the\\emulated device}};
  \draw[io] (gd.south) -- (em.north);
  \node[lab, text=ln-accent, anchor=west] at (19,34.3) {\iflangja{トラップ}{trap}};
  \draw[io, shorten >=0pt] (dd.south) -- (18,6);
  \end{scope}
  \node[pcap] at (59,-2) {\iflangja{(b) ハイパーバイザ直接}{(b) hypervisor-direct}};
  % ================= (c) split-device driver
  \begin{scope}[xshift=82mm]
  \draw[dev] (0,0) rectangle (36,6) node[midway] {\iflangja{デバイス}{device}};
  \draw[band] (0,8.5) rectangle (36,29);
  \node[hdr, anchor=west] at (1.2,26.5) {\iflangja{ハイパーバイザ}{hypervisor}};
  \draw[vm] (0,31) rectangle (17,55);
  \node[hdr] at (8.5,52.5) {\iflangja{ゲスト VM}{guest VM}};
  \node[box, fill=ln-accent!12, minimum width=14mm, minimum height=8mm] (fe)
    at (8.5,44) {\iflangja{フロント\\エンド}{front-end\\driver}};
  \draw[svm] (19,31) rectangle (36,55);
  \node[hdr] at (27.5,52.5) {\iflangja{サービス VM}{service VM}};
  \node[box, fill=ln-accent!12, minimum width=14mm, minimum height=8mm] (be)
    at (27.5,44) {\iflangja{バック\\エンド}{back-end\\driver}};
  \node[box, minimum width=14mm, minimum height=6mm] (nd) at (27.5,35.5)
    {\iflangja{ドライバ}{driver}};
  \draw[<->, thick, ln-accent] (fe.east) -- (be.west);
  \draw[->] (be.south) -- (nd.north);
  \draw[io] (nd.south) -- (27.5,6);
  \end{scope}
  \node[pcap] at (100,-2) {\iflangja{(c) 分割デバイスドライバ}{(c) split-device driver}};
\end{tikzpicture}
   (width ~118 mm)
\begin{tikzpicture}[x=1mm, y=1mm, font=\scriptsize\sffamily,
  >={Stealth[length=1.5mm]},
  dev/.style={draw, fill=white},
  band/.style={draw, fill=ln-accent!22},
  vm/.style={draw=ln-rule, rounded corners=2pt, fill=ln-box},
  svm/.style={draw=ln-accent, rounded corners=2pt, fill=ln-accent!8},
  gos/.style={draw, fill=ln-accent!12},
  box/.style={draw, fill=white, align=center, inner sep=0.5pt,
              font=\scriptsize\sffamily},
  hdr/.style={font=\scriptsize\sffamily\bfseries},
  lab/.style={font=\scriptsize\sffamily, text=ln-muted},
  io/.style={->, thick, ln-accent},
  pcap/.style={font=\footnotesize\sffamily, anchor=north}]
  % ================= (a) passthrough
  \draw[dev] (0,0) rectangle (36,6) node[midway] {\iflangja{デバイス}{device}};
  \draw[band] (0,8.5) rectangle (22,29);
  \node[hdr] at (11,26.5) {\iflangja{ハイパーバイザ}{hypervisor}};
  \node[box, minimum width=18mm, minimum height=8mm] (vd) at (11,16.5)
    {\iflangja{VMM レベル\\ドライバ}{VMM-level\\driver}};
  \draw[->, densely dashed] (vd.south) -- (11,6);
  \draw[vm] (0,31) rectangle (36,55);
  \node[hdr] at (18,52.5) {\iflangja{ゲスト VM}{guest VM}};
  \draw[gos] (2,33) rectangle (34,49.5);
  \node[lab, text=black] at (18,47.2) {\iflangja{ゲスト OS}{guest OS}};
  \node[box, minimum width=24mm, minimum height=7mm] at (18,39.5)
    {\iflangja{デバイスドライバ}{device driver}};
  \draw[io] (29,36) -- (29,6);
  \node[lab, rotate=90, text=ln-accent] at (32.2,19) {\iflangja{直接アクセス}{direct access}};
  \node[pcap] at (18,-2) {\iflangja{(a) パススルー}{(a) passthrough}};
  % ================= (b) hypervisor-direct
  \begin{scope}[xshift=41mm]
  \draw[dev] (0,0) rectangle (36,6) node[midway] {\iflangja{デバイス}{device}};
  \draw[band] (0,8.5) rectangle (36,29);
  \node[hdr, anchor=west] at (1.2,26.5) {\iflangja{ハイパーバイザ}{hypervisor}};
  \node[box, minimum width=30mm, minimum height=4.3mm] (em) at (18,21.7)
    {\iflangja{デバイスエミュレーション}{device emulation}};
  \node[box, minimum width=30mm, minimum height=4.3mm] (st) at (18,16.6)
    {\iflangja{I/O スタック}{I/O stack}};
  \node[box, minimum width=30mm, minimum height=4.3mm] (dd) at (18,11.5)
    {\iflangja{デバイスドライバ}{device driver}};
  \draw[vm] (0,31) rectangle (36,55);
  \node[hdr] at (18,52.5) {\iflangja{ゲスト VM}{guest VM}};
  \draw[gos] (2,33) rectangle (34,49.5);
  \node[lab, text=black] at (18,47.2) {\iflangja{ゲスト OS}{guest OS}};
  \node[box, minimum width=26mm, minimum height=7mm] (gd) at (18,39.5)
    {\iflangja{エミュレートされた\\デバイスのドライバ}{driver for the\\emulated device}};
  \draw[io] (gd.south) -- (em.north);
  \node[lab, text=ln-accent, anchor=west] at (19,34.3) {\iflangja{トラップ}{trap}};
  \draw[io, shorten >=0pt] (dd.south) -- (18,6);
  \end{scope}
  \node[pcap] at (59,-2) {\iflangja{(b) ハイパーバイザ直接}{(b) hypervisor-direct}};
  % ================= (c) split-device driver
  \begin{scope}[xshift=82mm]
  \draw[dev] (0,0) rectangle (36,6) node[midway] {\iflangja{デバイス}{device}};
  \draw[band] (0,8.5) rectangle (36,29);
  \node[hdr, anchor=west] at (1.2,26.5) {\iflangja{ハイパーバイザ}{hypervisor}};
  \draw[vm] (0,31) rectangle (17,55);
  \node[hdr] at (8.5,52.5) {\iflangja{ゲスト VM}{guest VM}};
  \node[box, fill=ln-accent!12, minimum width=14mm, minimum height=8mm] (fe)
    at (8.5,44) {\iflangja{フロント\\エンド}{front-end\\driver}};
  \draw[svm] (19,31) rectangle (36,55);
  \node[hdr] at (27.5,52.5) {\iflangja{サービス VM}{service VM}};
  \node[box, fill=ln-accent!12, minimum width=14mm, minimum height=8mm] (be)
    at (27.5,44) {\iflangja{バック\\エンド}{back-end\\driver}};
  \node[box, minimum width=14mm, minimum height=6mm] (nd) at (27.5,35.5)
    {\iflangja{ドライバ}{driver}};
  \draw[<->, thick, ln-accent] (fe.east) -- (be.west);
  \draw[->] (be.south) -- (nd.north);
  \draw[io] (nd.south) -- (27.5,6);
  \end{scope}
  \node[pcap] at (100,-2) {\iflangja{(c) 分割デバイスドライバ}{(c) split-device driver}};
\end{tikzpicture}
   (width ~118 mm)
\begin{tikzpicture}[x=1mm, y=1mm, font=\scriptsize\sffamily,
  >={Stealth[length=1.5mm]},
  dev/.style={draw, fill=white},
  band/.style={draw, fill=ln-accent!22},
  vm/.style={draw=ln-rule, rounded corners=2pt, fill=ln-box},
  svm/.style={draw=ln-accent, rounded corners=2pt, fill=ln-accent!8},
  gos/.style={draw, fill=ln-accent!12},
  box/.style={draw, fill=white, align=center, inner sep=0.5pt,
              font=\scriptsize\sffamily},
  hdr/.style={font=\scriptsize\sffamily\bfseries},
  lab/.style={font=\scriptsize\sffamily, text=ln-muted},
  io/.style={->, thick, ln-accent},
  pcap/.style={font=\footnotesize\sffamily, anchor=north}]
  % ================= (a) passthrough
  \draw[dev] (0,0) rectangle (36,6) node[midway] {\iflangja{デバイス}{device}};
  \draw[band] (0,8.5) rectangle (22,29);
  \node[hdr] at (11,26.5) {\iflangja{ハイパーバイザ}{hypervisor}};
  \node[box, minimum width=18mm, minimum height=8mm] (vd) at (11,16.5)
    {\iflangja{VMM レベル\\ドライバ}{VMM-level\\driver}};
  \draw[->, densely dashed] (vd.south) -- (11,6);
  \draw[vm] (0,31) rectangle (36,55);
  \node[hdr] at (18,52.5) {\iflangja{ゲスト VM}{guest VM}};
  \draw[gos] (2,33) rectangle (34,49.5);
  \node[lab, text=black] at (18,47.2) {\iflangja{ゲスト OS}{guest OS}};
  \node[box, minimum width=24mm, minimum height=7mm] at (18,39.5)
    {\iflangja{デバイスドライバ}{device driver}};
  \draw[io] (29,36) -- (29,6);
  \node[lab, rotate=90, text=ln-accent] at (32.2,19) {\iflangja{直接アクセス}{direct access}};
  \node[pcap] at (18,-2) {\iflangja{(a) パススルー}{(a) passthrough}};
  % ================= (b) hypervisor-direct
  \begin{scope}[xshift=41mm]
  \draw[dev] (0,0) rectangle (36,6) node[midway] {\iflangja{デバイス}{device}};
  \draw[band] (0,8.5) rectangle (36,29);
  \node[hdr, anchor=west] at (1.2,26.5) {\iflangja{ハイパーバイザ}{hypervisor}};
  \node[box, minimum width=30mm, minimum height=4.3mm] (em) at (18,21.7)
    {\iflangja{デバイスエミュレーション}{device emulation}};
  \node[box, minimum width=30mm, minimum height=4.3mm] (st) at (18,16.6)
    {\iflangja{I/O スタック}{I/O stack}};
  \node[box, minimum width=30mm, minimum height=4.3mm] (dd) at (18,11.5)
    {\iflangja{デバイスドライバ}{device driver}};
  \draw[vm] (0,31) rectangle (36,55);
  \node[hdr] at (18,52.5) {\iflangja{ゲスト VM}{guest VM}};
  \draw[gos] (2,33) rectangle (34,49.5);
  \node[lab, text=black] at (18,47.2) {\iflangja{ゲスト OS}{guest OS}};
  \node[box, minimum width=26mm, minimum height=7mm] (gd) at (18,39.5)
    {\iflangja{エミュレートされた\\デバイスのドライバ}{driver for the\\emulated device}};
  \draw[io] (gd.south) -- (em.north);
  \node[lab, text=ln-accent, anchor=west] at (19,34.3) {\iflangja{トラップ}{trap}};
  \draw[io, shorten >=0pt] (dd.south) -- (18,6);
  \end{scope}
  \node[pcap] at (59,-2) {\iflangja{(b) ハイパーバイザ直接}{(b) hypervisor-direct}};
  % ================= (c) split-device driver
  \begin{scope}[xshift=82mm]
  \draw[dev] (0,0) rectangle (36,6) node[midway] {\iflangja{デバイス}{device}};
  \draw[band] (0,8.5) rectangle (36,29);
  \node[hdr, anchor=west] at (1.2,26.5) {\iflangja{ハイパーバイザ}{hypervisor}};
  \draw[vm] (0,31) rectangle (17,55);
  \node[hdr] at (8.5,52.5) {\iflangja{ゲスト VM}{guest VM}};
  \node[box, fill=ln-accent!12, minimum width=14mm, minimum height=8mm] (fe)
    at (8.5,44) {\iflangja{フロント\\エンド}{front-end\\driver}};
  \draw[svm] (19,31) rectangle (36,55);
  \node[hdr] at (27.5,52.5) {\iflangja{サービス VM}{service VM}};
  \node[box, fill=ln-accent!12, minimum width=14mm, minimum height=8mm] (be)
    at (27.5,44) {\iflangja{バック\\エンド}{back-end\\driver}};
  \node[box, minimum width=14mm, minimum height=6mm] (nd) at (27.5,35.5)
    {\iflangja{ドライバ}{driver}};
  \draw[<->, thick, ln-accent] (fe.east) -- (be.west);
  \draw[->] (be.south) -- (nd.north);
  \draw[io] (nd.south) -- (27.5,6);
  \end{scope}
  \node[pcap] at (100,-2) {\iflangja{(c) 分割デバイスドライバ}{(c) split-device driver}};
\end{tikzpicture}

  \caption{デバイスの仮想化のモデル．(a) VM 自身のドライバがデバイスに直接
    アクセスする．(b) ハイパーバイザがデバイスアクセスを横取りし，自分の I/O
    スタックとドライバでそれをエミュレートする．(c) ゲスト内のフロントエンド
    ドライバが，サービス VM 内のバックエンドドライバに要求を渡し，バックエンド
    ドライバがデバイスを駆動する．}
  \label{fig:l12-devices}
\end{figure}

\needspace{6\baselineskip}
\subsection{パススルーモデル}

\begin{definition}[パススルーモデル]
  \term[ぱすするーもでる]{パススルーモデル}[passthrough model]では，VMM
  レベルのドライバがデバイスのアクセス権を設定し，VM が自分のドライバで，
  ハイパーバイザを迂回してデバイスに直接アクセスする．
\end{definition}

VM はデバイスへの排他的なアクセスを得て，そのデバイス操作は物理マシン上と同じ
経路をたどる．この \emph{VMM バイパス}は，第11章の OS
バイパス\index{OSばいぱす@OSバイパス}に対応するもので，最良の性能を与える．
このモデルには3つの欠点がある．
\begin{itemize}
  \item デバイスの共有が難しい．ある VM から別の VM へ割り当て直さなければ
    ならない．\margin{SR-IOV（PCI-SIG，2007 年）はこれを和らげる．1 枚の
    カードが1つの物理ファンクションと複数の仮想ファンクション---仕様上は
    256 個まで，実際には数十個---を提示し，その各々を VM に割り当てられる．}
  \item VM 内のドライバがハードウェアを直接操作するので，マシンはそのドライバ
    が期待する種類のデバイスをちょうど備えていなければならない．
  \item VM の状態の一部がデバイスの中にあり，移送先のマシンがその状態を移せる
    同一のデバイスを備えていなければならないので，マイグレーションが厄介で
    ある．
\end{itemize}

\subsection{ハイパーバイザ直接モデル}

\begin{definition}[ハイパーバイザ直接モデル]
  \term[はいぱーばいざちょくせつもでる]{ハイパーバイザ直接モデル}%
  [hypervisor-direct model]では，ハイパーバイザが VM のデバイスアクセスを
  すべて横取りし，デバイス操作をエミュレートする．すなわち，その操作を汎用の
  I/O 操作に翻訳し，ハイパーバイザ内にある I/O スタックを通し，物理デバイス用の
  ハイパーバイザ自身のドライバを呼び出す．
\end{definition}

主な利点は，VM が物理デバイスから切り離されることである．\margin{エミュレート
されるデバイスはわざと古く，ありふれたものが選ばれる．QEMU の PC マシンは
Intel e1000 のネットワークカードと IDE のディスクコントローラを提示する．
どのゲスト OS もそのドライバをすでに持っているからである．}VM はエミュレート
されたデバイスを見る．それは物理デバイスとは無関係に，どのマシンでも同じで
ありうる．デバイスは VM の間で共有でき，マイグレーションは簡単になり，各
デバイスの固有事情はハイパーバイザが扱う．VMware ESX はこのモデルを用いる．
代価は，横取りとエミュレーションがデバイス操作のたびに加えるレイテンシと，
ハイパーバイザ内部のデバイスドライバのエコシステムの複雑さである．ハイパー
バイザは，サポートするすべてのデバイスのドライバを含み，しかも堅牢であり続け
なければならない．

\subsection{分割デバイスドライバモデル}

\begin{definition}[分割デバイスドライバモデル]
  \term[ぶんかつでばいすどらいばもでる]{分割デバイスドライバモデル}%
  [split-device driver model]では，デバイスへのアクセスを，協調する2つの
  ドライバが制御する．1つはゲスト VM 内の
  \term[ふろんとえんどどらいば]{フロントエンドドライバ}[front-end driver]で，
  ゲスト OS にデバイスインタフェースを提示する．もう1つはサービス VM または
  ホスト OS 内の
  \term[ばっくえんどどらいば]{バックエンドドライバ}[back-end driver]で，
  通常のドライバを通じて物理デバイスにアクセスする．
\end{definition}

フロントエンドドライバは，トラップしてエミュレートしなければならないような
デバイスレジスタを操作しない．送信すべきパケットのような完結した要求を
バックエンドドライバに渡す．Xen では，2つの VM の間で共有されるメモリ内の
記述子リングを通じて渡す．\margin{Linux の virtio（Russell，2008 年．
2016 年以降は OASIS の標準）が，共通のフロントエンド・バックエンドの
インタフェースである．その virtqueue は第9章の共有メモリである．}%
これによってハイパーバイザ直接モデルのエミュレーションのオーバーヘッドが
なくなり，あるデバイスへの要求はすべてバックエンド
ドライバを通るので，サービス VM はそのデバイスが VM の間でどう共有されるかを
管理できる．代価はゲストの改変である．そのドライバを，ハイパーバイザ向けに
書かれたフロントエンドドライバに置き換えなければならないので，このモデルは
準仮想化されたゲストと，そうしたドライバが導入されたゲストに適用される．
\cref{tab:l12-devices}に3つのモデルをまとめる．

\begin{table}[htbp]
  \centering
  \caption{デバイスの仮想化のモデルの比較．}
  \label{tab:l12-devices}
  \small
  \begin{tabular}{@{}>{\raggedright}p{18mm}>{\raggedright}p{19mm}%
      >{\raggedright}p{20mm}>{\raggedright}p{23mm}%
      >{\raggedright\arraybackslash}p{23mm}@{}}
    \toprule
    モデル & 物理デバイスを駆動するもの & ゲスト内のドライバ & 共有と
      マイグレーション & デバイス操作ごとの代価 \\
    \midrule
    パススルー & ゲスト VM & ネイティブドライバ & 難しい & ネイティブ以上の
      代価なし \\
    \addlinespace
    ハイパーバイザ直接 & ハイパーバイザ & エミュレートされたデバイスの
      ドライバ & 可能 & 横取りとエミュレーション \\
    \addlinespace
    分割デバイスドライバ & サービス VM またはホスト OS & 改変された
      フロントエンド & 可能，バックエンドが管理 & バックエンドへの要求の
      受け渡し \\
    \bottomrule
  \end{tabular}
\end{table}

\begin{keypoint}
  デバイスには標準のインタフェースがないので，その仮想化は性能と柔軟性を
  引き換えにする．パススルーは最速だが VM をデバイスに縛りつけ，ハイパー
  バイザ直接はエミュレーションの代価で VM を切り離し，分割デバイスドライバ
  モデルはゲストのドライバを改変することでエミュレーションを避ける．
\end{keypoint}

\section{仮想化のためのハードウェア支援}
\label{sec:l12-hw}

2005 年頃，AMD と Intel はプロセッサに仮想化のための支援を加えた．AMD は
Pacifica というコード名の技術，現在の AMD-V であり，Intel は Vanderpool，
現在の Intel VT である．\source{P3L6，ハードウェアによる仮想化，x86 の VT
革命；Uhlig ら \cite{uhlig2005}；Rosenblum と Garfinkel
\cite{rosenblum2005}．}

\begin{definition}[ハードウェア支援仮想化]
  プロセッサ自体が，それを仮想化するために必要な機構を提供し，改変されていない
  ゲストオペレーティングシステムが，バイナリ変換も準仮想化もなしにトラップ
  アンドエミュレートで動作するような仮想化を，
  \term[はーどうぇあしえんかそうか]{ハードウェア支援仮想化}%
  [hardware-assisted virtualization]と呼ぶ．
\end{definition}

これらの拡張は，これまでの節で述べた問題に対処する．
\begin{description}
  \item[x86 ISA の穴をふさぐ] 非ルートモードではゲスト OS はリング 0 で動作
    し，そこでは \texttt{POPF} のような命令もゲストの期待どおりに振る舞う．
    また，センシティブ命令，制御レジスタへのアクセス，I/O，割り込みなど，
    ハイパーバイザが制御しなければならない操作はすべて，VM exit を起こすように
    できる．これによってこのアーキテクチャは，トラップアンドエミュレートで
    仮想化可能になる．
  \item[ルートモードと非ルートモード] \cref{sec:l12-rings}のモードは，AMD では
    ホストモードとゲストモードと呼ばれ，ハイパーバイザが制御を保ったまま
    ゲスト OS をリング 0 で動作させる．
  \item[VM 制御構造] VM ごと，より正確には仮想 CPU ごとに1つのデータ構造で，
    ハードウェアがそれを読み書きする．仮想 CPU の状態を保持し，VM exit で保存
    され VM entry で復元される．また，どの命令とイベントが VM exit を起こすかを
    指定するので，ハイパーバイザは VM ごとに何をトラップするかを決められる．
  \item[拡張ページテーブル] MMU が\cref{fig:l12-memory}の2段階の変換を自ら
    行う．ゲストのページテーブルは従来どおり仮想アドレスを物理アドレスへ対応
    付け，ハイパーバイザが維持する
    \term[かくちょうぺーじてーぶる]{拡張ページテーブル}[extended page table]%
    （EPT）\index{EPT|see{拡張ページテーブル}}が物理アドレスをマシンアドレスへ
    対応付ける．TLB ミスの際，MMU は両方のテーブルを辿り，TLB は結果として
    得られる仮想アドレスからマシンアドレスへの変換をキャッシュする．%
    \tip{ゲストの $g$ 段の下に拡張の $n$ 段があるとき，1回の探索には
    $(g+1)(n+1)-1$ 回の参照が要る．$g=n=4$ なら 24 回であり，ネイティブの
    4 回に対する代価である（AMD-V ネステッドページングのホワイトペーパー）．}%
    ゲスト OS は VM exit なしに自分のページテーブルを変更し切り替えられる
    ようになり，
    シャドウページテーブルは不要になる．AMD における同等のものは，ネステッド
    ページテーブルと呼ばれる．
  \item[タグ付き TLB] \term[たぐつきTLB]{タグ付きTLB}[tagged TLB]は，各
    エントリに，それが属する VM の識別子を付ける．異なる VM のエントリが TLB の
    中に共存できるので，VM の切り替えでそれをフラッシュする必要がない．
  \item[マルチキューデバイスと割り込みのルーティング] ネットワークカードの
    ようなデバイスは，複数の論理インタフェース，すなわちキューを提供し，その
    各々を別々の VM に割り当てられる．割り込みは，そのキューを所有する VM が
    動作しているコアへ配送される．
  \item[セキュリティと管理の支援] さらなる拡張が，ハイパーバイザを保護し，VM
    の管理を支える．
\end{description}
新しい命令が，非ルートモードへの遷移，VM 制御構造の読み書き，ハイパーコールの
発行を行う．\margin{Intel では，VM entry に \texttt{VMLAUNCH} と
\texttt{VMRESUME}，制御構造に \texttt{VMREAD} と \texttt{VMWRITE}，
ハイパーコールに \texttt{VMCALL} を用いる．}KVM のようなホスト型ハイパー
バイザは，これらに依存する．

Intel は支援をプロセッサから外側へ向かって段階的に広げた．\emph{VT-x} は
プロセッサを対象とし，ルートモードと非ルートモード，VM exit と VM entry，VM
制御構造を，そして後のプロセッサでは拡張ページテーブルとタグ付き TLB エントリ
を含む．\margin{VT-x は 2005 年 11 月の Pentium~4 662 と 672 で，AMD-V は
2006 年 5 月に出荷された．ページテーブルと TLB の支援はさらに後で，ネステッド
ページテーブルと ASID タグが Barcelona（2007 年），拡張ページテーブルと
16 ビットの VPID タグが Nehalem（2008 年）である．}%
\emph{VT-d}（directed I/O）はチップセットに I/O メモリ管理ユニットを加え，
DMA（第11章）で使われるアドレスを変換し割り込みを再対応付けするので，
パススルーモデルで VM に割り当てられたデバイスが，他の VM のメモリにアクセス
できないようにする．\emph{VT-c}（connectivity）はネットワークデバイスに仮想化
支援をもたらす．デバイスごとに複数のキューを備え，自らを複数の仮想デバイスと
して提示して，その各々を VM に割り当てられるようにする．

\begin{table}[htbp]
  \centering
  \caption{x86 におけるプロセッサの仮想化の手法．}
  \label{tab:l12-cpu}
  \small
  \begin{tabular}{@{}>{\raggedright}p{24mm}>{\raggedright}p{16mm}%
      >{\raggedright}p{27mm}>{\raggedright}p{22mm}%
      >{\raggedright\arraybackslash}p{14mm}@{}}
    \toprule
    手法 & ゲスト OS & 問題のある命令 & 主なオーバーヘッド & 例 \\
    \midrule
    トラップアンドエミュレートのみ，2005 年以前 & 改変なし & 一部が黙って失敗
      する: 不正確 & トラップ & --- \\
    \addlinespace
    バイナリ変換 & 改変なし & 実行時に書き換え & 変換，エミュレーション
      & VMware \\
    \addlinespace
    準仮想化 & 改変あり & ハイパーコールに置換 & ハイパーコール & Xen \\
    \addlinespace
    ハードウェア支援 & 改変なし & VM exit を起こす & VM exit & KVM \\
    \bottomrule
  \end{tabular}
\end{table}

\begin{keypoint}[章のまとめ]
  仮想化は1台のマシン上で複数のオペレーティングシステムを並行に動かす．
  ハイパーバイザは各 VM に，効率的で隔離されたマシンの複製を与える．ベアメタル
  型ハイパーバイザはハードウェア上で動作し，通常はドライバを所有するサービス
  VM を伴う．ホスト型ハイパーバイザはホスト OS を拡張する．CPU はトラップ
  アンドエミュレートで仮想化されるが，2005 年以前の x86 プロセッサはこれを
  完全にはサポートしていなかったため，改変しないゲストにはバイナリ変換が，
  改変するゲストにはハイパーコールを用いる準仮想化が使われた．メモリの仮想化は
  物理アドレスの下にマシンアドレスを加え，シャドウページテーブルか，検証され
  まとめられたページテーブル更新によって扱われる．デバイスはパススルー，
  ハイパーバイザ直接，分割デバイスドライバの各モデルで仮想化される．ルート
  モードと非ルートモード，VM exit，拡張ページテーブル，タグ付き TLB を備えた
  ハードウェア支援仮想化は，これらのオーバーヘッドの大半を取り除く．
\end{keypoint}

\end{document}
